\documentclass[../main.tex]{subfiles}
\begin{document}

\chapter{Distributed Shared Memory}
\lessoninfo{
  video={P4L3},
  textbook={Nitzberg and Lo \cite{nitzberg1991}; Tanenbaum and Bos \cite{tanenbaum2015} ch.~8; Tanenbaum and van Steen \cite{tanenbaum2007} ch.~7},
  keywords={distributed shared memory, RDMA, sharing granularity, false sharing, access algorithm, migration, replication, consistency management, page fault, consistency model, strict consistency, sequential consistency, causal consistency, weak consistency},
}

Lesson~14 studied distributed file systems, in which servers own the files
and serve them to clients that cache parts of them.  In many distributed
applications, however, every node both owns a part of the state and serves
it to the others.  This lesson examines distributed shared memory, which
lets processes on different machines share memory as if they ran on a
single shared memory multiprocessor, as an example of how such distributed
state is managed.  It discusses the design choices of a DSM
system---granularity, access patterns, migration versus replication and
consistency management---how a DSM locates its pages and uses the MMU to
intercept accesses, and the consistency models that define what processes
observe when they share memory.

\section{Peer distributed applications}
\label{sec:l15-peer}

In the distributed file systems of Lesson~14 the roles of the machines are
clearly divided.\source{P4L3, lesson preview, reviewing distributed file
systems, peer distributed applications; Tanenbaum and van Steen
\cite{tanenbaum2007} ch.~2.}  Servers own and manage the state, the files,
and provide the service, file access; clients
send requests to the servers.  Clients cache parts of the state, which
lowers the latency of their accesses and lets a server support more
clients, but the state itself belongs to the servers.

In many distributed applications no such division exists.  In big data
analytics, web search or content sharing,\margin{Lesson~16 examines such
applications in the datacenter, where a framework sends each computation to
the node that already holds its input.} the state of the application is
distributed over all of its nodes: every node owns a portion of the state
and provides service on that portion to the other nodes, while it also
uses the state owned by others.  Every node is thus both a client and a
server, and all nodes are peers; we call such applications \emph{peer
distributed applications}.  The word \enquote{peer} is used loosely here.
In a peer-to-peer system (Lesson~14)\index{peer-to-peer} even the control
and management of the overall system are carried out by all peers, whereas
a peer distributed application may assign management tasks, such as
keeping track of which node holds which part of the state, to designated
nodes.

Managing state that is spread over peers raises three issues, which recur
throughout this lesson:\margin{The first two arise inside one machine as
well: NUMA-aware scheduling (Lesson~7) moves the computation to its memory
instead.}
\begin{description}
  \item[Placement] The node on which each part of the state resides.
    Access is fastest when the state resides on the node of the processes
    that access it most.
  \item[Migration] How and when state that resides on another node is
    brought to the node of the process that accesses it, by moving it or by
    copying it.\caution{The migration of this lesson moves \emph{memory}
    between nodes.  It is not the process migration of Lesson~8, which
    transfers a whole process to another machine.}
  \item[Sharing rules] The rules that determine the order in which
    concurrent operations on the same state take effect, and when each node
    observes the updates made by the others.
\end{description}
Distributed shared memory, a peer distributed application whose state is
memory, is the example on which the lesson studies these issues.

\section{Distributed shared memory}
\label{sec:l15-dsm}

A shared memory multiprocessor (Lesson~7)\index{shared memory
multiprocessor} lets threads on different CPUs cooperate through memory
that all of them access, but its memory is limited to what one machine can
hold.\source{P4L3, distributed shared memory; Nitzberg and Lo
\cite{nitzberg1991}; Tanenbaum and Bos ch.~8.}  Distributed
shared memory extends the model to many machines.

\begin{definition}[Distributed shared memory]
  A service that manages memory across the nodes of a network, so that the
  applications running on top of it have the illusion of running on a
  single shared memory machine, is called
  \term{distributed shared memory}[分散共有メモリ]%
  \index{DSM|see{distributed shared memory}} (DSM).  Processes on any node
  read and write one shared address space, whose contents reside in the
  physical memories of the different nodes.
\end{definition}

Like every peer distributed application, each node of a DSM system
(\cref{fig:l15-dsm})
\begin{itemize}
  \item owns state: a portion of the physical memory of the system, namely
    the part of its own memory that it contributes to the DSM;
  \item provides service on that state: read and write operations on its
    memory, which may be issued by processes on any node;
  \item takes part in consistency protocols with the other nodes, which
    ensure that accesses to the shared memory remain meaningful when
    several nodes access the same memory---for example, that a read returns
    the value of the most recent write.\caution{This is not the shared
    memory IPC of Lesson~9: there processes share the memory of one machine,
    and the hardware keeps their views consistent.}
\end{itemize}

\begin{figure}[htbp]
  \centering
  % Distributed shared memory: processes on three nodes read and write one
% shared address space; every page of it resides in the memory that one of
% the nodes contributes, and remote accesses cross the network.
% Usage: % Distributed shared memory: processes on three nodes read and write one
% shared address space; every page of it resides in the memory that one of
% the nodes contributes, and remote accesses cross the network.
% Usage: % Distributed shared memory: processes on three nodes read and write one
% shared address space; every page of it resides in the memory that one of
% the nodes contributes, and remote accesses cross the network.
% Usage: \input{figures/l15-dsm}   (width ~118 mm)
\begin{tikzpicture}[x=1mm, y=1mm, font=\scriptsize\sffamily,
  >={Stealth[length=1.3mm]},
  mach/.style={draw, rounded corners=0.8mm, fill=white},
  proc/.style={draw, circle, fill=white, inner sep=0pt, minimum size=4.8mm},
  dsm/.style={draw, fill=ln-box, minimum width=30mm, minimum height=4.4mm, inner sep=0pt},
  frm/.style={draw, minimum width=8.4mm, minimum height=4.6mm, inner sep=0pt},
  cell/.style={draw, minimum width=13mm, minimum height=5mm, inner sep=0pt},
  lab/.style={font=\scriptsize\sffamily, text=ln-muted, inner sep=0.8pt},
  who/.style={font=\scriptsize\sffamily\bfseries, anchor=north west, inner sep=1pt}]
  \colorlet{lnNodeA}{ln-accent!22}
  \colorlet{lnNodeB}{ln-tip!22}
  \colorlet{lnNodeC}{ln-caution!16}
  % shared address space
  \node[lab, anchor=south west] at (1,40.8)
    {\iflangja{共有アドレス空間（すべてのプロセスから見える）}{shared address space (seen by every process)}};
  \foreach \ci/\cc in {0/lnNodeA,1/lnNodeB,2/lnNodeC,3/lnNodeA,4/lnNodeB,5/lnNodeC,6/lnNodeA,7/lnNodeB,8/lnNodeC} {
    \node[cell, fill=\cc] at ({7.5+13*\ci},38) {$p_{\ci}$};
  }
  % nodes
  \foreach \nk/\nx/\ncol/\npa/\npb/\fa/\fb/\fc in
      {1/0/lnNodeA/1/2/0/3/6, 2/41/lnNodeB/3/4/1/4/7, 3/82/lnNodeC/5/6/2/5/8} {
    \draw[mach] (\nx,0) rectangle ++(36,26);
    \node[who] at (\nx+0.5,25.5) {\iflangja{ノード}{node} \nk};
    \node[proc] at (\nx+15,19.5) {P\npa};
    \node[proc] at (\nx+27,19.5) {P\npb};
    \draw[<->] (\nx+15,21.9) -- (\nx+15,35.4);
    \draw[<->] (\nx+27,21.9) -- (\nx+27,35.4);
    \node[dsm] at (\nx+18,12) {DSM};
    \foreach \fj/\fp in {0/\fa,1/\fb,2/\fc} {
      \node[frm, fill=\ncol] at (\nx+8.5+9.5*\fj,4.6) {$p_{\fp}$};
    }
    \draw[thick, ln-muted] (\nx+18,0) -- (\nx+18,-5);
  }
  \node[lab, anchor=west] at (27.6,30.5) {\iflangja{読み書き}{read / write}};
  % network
  \draw[line width=1.2pt, ln-muted] (0,-5) -- (118,-5);
  \node[lab, anchor=north] at (59,-5.8)
    {\iflangja{ネットワーク（RDMA 対応のインターコネクトなど）}{network (e.g.\ an interconnect with RDMA)}};
\end{tikzpicture}
   (width ~118 mm)
\begin{tikzpicture}[x=1mm, y=1mm, font=\scriptsize\sffamily,
  >={Stealth[length=1.3mm]},
  mach/.style={draw, rounded corners=0.8mm, fill=white},
  proc/.style={draw, circle, fill=white, inner sep=0pt, minimum size=4.8mm},
  dsm/.style={draw, fill=ln-box, minimum width=30mm, minimum height=4.4mm, inner sep=0pt},
  frm/.style={draw, minimum width=8.4mm, minimum height=4.6mm, inner sep=0pt},
  cell/.style={draw, minimum width=13mm, minimum height=5mm, inner sep=0pt},
  lab/.style={font=\scriptsize\sffamily, text=ln-muted, inner sep=0.8pt},
  who/.style={font=\scriptsize\sffamily\bfseries, anchor=north west, inner sep=1pt}]
  \colorlet{lnNodeA}{ln-accent!22}
  \colorlet{lnNodeB}{ln-tip!22}
  \colorlet{lnNodeC}{ln-caution!16}
  % shared address space
  \node[lab, anchor=south west] at (1,40.8)
    {\iflangja{共有アドレス空間（すべてのプロセスから見える）}{shared address space (seen by every process)}};
  \foreach \ci/\cc in {0/lnNodeA,1/lnNodeB,2/lnNodeC,3/lnNodeA,4/lnNodeB,5/lnNodeC,6/lnNodeA,7/lnNodeB,8/lnNodeC} {
    \node[cell, fill=\cc] at ({7.5+13*\ci},38) {$p_{\ci}$};
  }
  % nodes
  \foreach \nk/\nx/\ncol/\npa/\npb/\fa/\fb/\fc in
      {1/0/lnNodeA/1/2/0/3/6, 2/41/lnNodeB/3/4/1/4/7, 3/82/lnNodeC/5/6/2/5/8} {
    \draw[mach] (\nx,0) rectangle ++(36,26);
    \node[who] at (\nx+0.5,25.5) {\iflangja{ノード}{node} \nk};
    \node[proc] at (\nx+15,19.5) {P\npa};
    \node[proc] at (\nx+27,19.5) {P\npb};
    \draw[<->] (\nx+15,21.9) -- (\nx+15,35.4);
    \draw[<->] (\nx+27,21.9) -- (\nx+27,35.4);
    \node[dsm] at (\nx+18,12) {DSM};
    \foreach \fj/\fp in {0/\fa,1/\fb,2/\fc} {
      \node[frm, fill=\ncol] at (\nx+8.5+9.5*\fj,4.6) {$p_{\fp}$};
    }
    \draw[thick, ln-muted] (\nx+18,0) -- (\nx+18,-5);
  }
  \node[lab, anchor=west] at (27.6,30.5) {\iflangja{読み書き}{read / write}};
  % network
  \draw[line width=1.2pt, ln-muted] (0,-5) -- (118,-5);
  \node[lab, anchor=north] at (59,-5.8)
    {\iflangja{ネットワーク（RDMA 対応のインターコネクトなど）}{network (e.g.\ an interconnect with RDMA)}};
\end{tikzpicture}
   (width ~118 mm)
\begin{tikzpicture}[x=1mm, y=1mm, font=\scriptsize\sffamily,
  >={Stealth[length=1.3mm]},
  mach/.style={draw, rounded corners=0.8mm, fill=white},
  proc/.style={draw, circle, fill=white, inner sep=0pt, minimum size=4.8mm},
  dsm/.style={draw, fill=ln-box, minimum width=30mm, minimum height=4.4mm, inner sep=0pt},
  frm/.style={draw, minimum width=8.4mm, minimum height=4.6mm, inner sep=0pt},
  cell/.style={draw, minimum width=13mm, minimum height=5mm, inner sep=0pt},
  lab/.style={font=\scriptsize\sffamily, text=ln-muted, inner sep=0.8pt},
  who/.style={font=\scriptsize\sffamily\bfseries, anchor=north west, inner sep=1pt}]
  \colorlet{lnNodeA}{ln-accent!22}
  \colorlet{lnNodeB}{ln-tip!22}
  \colorlet{lnNodeC}{ln-caution!16}
  % shared address space
  \node[lab, anchor=south west] at (1,40.8)
    {\iflangja{共有アドレス空間（すべてのプロセスから見える）}{shared address space (seen by every process)}};
  \foreach \ci/\cc in {0/lnNodeA,1/lnNodeB,2/lnNodeC,3/lnNodeA,4/lnNodeB,5/lnNodeC,6/lnNodeA,7/lnNodeB,8/lnNodeC} {
    \node[cell, fill=\cc] at ({7.5+13*\ci},38) {$p_{\ci}$};
  }
  % nodes
  \foreach \nk/\nx/\ncol/\npa/\npb/\fa/\fb/\fc in
      {1/0/lnNodeA/1/2/0/3/6, 2/41/lnNodeB/3/4/1/4/7, 3/82/lnNodeC/5/6/2/5/8} {
    \draw[mach] (\nx,0) rectangle ++(36,26);
    \node[who] at (\nx+0.5,25.5) {\iflangja{ノード}{node} \nk};
    \node[proc] at (\nx+15,19.5) {P\npa};
    \node[proc] at (\nx+27,19.5) {P\npb};
    \draw[<->] (\nx+15,21.9) -- (\nx+15,35.4);
    \draw[<->] (\nx+27,21.9) -- (\nx+27,35.4);
    \node[dsm] at (\nx+18,12) {DSM};
    \foreach \fj/\fp in {0/\fa,1/\fb,2/\fc} {
      \node[frm, fill=\ncol] at (\nx+8.5+9.5*\fj,4.6) {$p_{\fp}$};
    }
    \draw[thick, ln-muted] (\nx+18,0) -- (\nx+18,-5);
  }
  \node[lab, anchor=west] at (27.6,30.5) {\iflangja{読み書き}{read / write}};
  % network
  \draw[line width=1.2pt, ln-muted] (0,-5) -- (118,-5);
  \node[lab, anchor=north] at (59,-5.8)
    {\iflangja{ネットワーク（RDMA 対応のインターコネクトなど）}{network (e.g.\ an interconnect with RDMA)}};
\end{tikzpicture}

  \caption{Distributed shared memory.  Every page $p_i$ of the shared
    address space resides in memory contributed by one of the nodes; the
    DSM layer of each node serves accesses to local pages and obtains
    remote pages over the network.}
  \label{fig:l15-dsm}
\end{figure}

DSM permits applications to scale beyond the memory limits of a single
machine: an application that needs more memory than one machine holds
obtains it from other machines.  It also provides shared memory at a lower
cost, because a single machine with a very large memory is
disproportionately more expensive than several commodity machines with
moderate memories.  The price is latency.\margin{An access to local memory
takes some \SI{100}{\nano\second} and a round trip over datacenter Ethernet
tens of microseconds, whereas a one-sided RDMA read over a
\SI{200}{\giga\bit\per\second} InfiniBand port takes about
\SI{1}{\micro\second} (NVIDIA ConnectX adapter specifications).}  An access
to memory that resides on another node crosses the network, which is orders
of magnitude slower than an access to local memory.  Commodity
interconnects reduce this cost by supporting remote direct memory access.

\begin{definition}[Remote direct memory access]
  With \term{remote direct memory access}[RDMA]%
  \index{RDMA|see{remote direct memory access}} (RDMA), a network interface
  transfers data directly between the memory of its own machine and the
  memory of a remote machine, without involving the CPU or the operating
  system of the remote machine in the transfer.
\end{definition}

RDMA extends direct memory access\index{direct memory access} (Lesson~11)
across the network.  It is typically combined with operating system
bypass\index{operating system bypass} (Lesson~11), so that the kernel of
neither machine lies on the data path.\margin{InfiniBand provides RDMA
natively; RoCE and iWARP are RDMA protocols carried over Ethernet.}

\section{Hardware and software DSM}
\label{sec:l15-hwsw}

DSM can be provided by hardware, by software, or by a combination of
both.\source{P4L3, hardware vs.\ software DSM; Nitzberg and Lo
\cite{nitzberg1991}.}
\begin{description}
  \item[Hardware-supported DSM] relies on a special-purpose interconnect.
    The operating system of each node is given the impression that it
    manages a physical memory much larger than the node's own, spanning
    several nodes.  An access to an address that resides on another node is
    translated by the network interface card (NIC) into a message to that
    node, and the NICs take part in every aspect of managing the shared
    memory, including atomic operations on remote memory, without involving
    the CPU of the remote node.  Such hardware is expensive; it is used in
    high-end systems where performance is paramount, as in high-performance
    computing.\margin{The Scalable Coherent Interface (IEEE 1596-1992),
    used by Dolphin's adapters and by the Sequent NUMA-Q, is such an
    interconnect: it keeps the memories of the nodes coherent in hardware.}
  \item[Software-supported DSM] performs everything in software: locating
    the node that holds the memory, fetching it, detecting accesses and
    keeping copies consistent.  The software is part of the operating
    system or of the runtime of a programming language; the first shares
    pages, the second objects (\cref{sec:l15-granularity}).
  \item[Hybrid DSM] divides the work between the two.  Hardware, the MMU in
    particular, performs what must happen on every access: it translates
    addresses and raises an exception when an access is not permitted by
    the current mapping.  Software makes the decisions that require
    knowledge of the system as a whole, for example which pages to
    prefetch\index{prefetching}, that is, to obtain from other nodes before
    they are accessed.
\end{description}

\section{Sharing granularity}
\label{sec:l15-granularity}

The first design decision of a DSM is the unit of sharing: the unit of
memory that the DSM transfers between nodes, replicates and keeps
consistent as a whole.\source{P4L3, DSM design: sharing granularity;
Nitzberg and Lo; Tanenbaum and Bos ch.~8.}  On a shared memory
multiprocessor the hardware keeps the CPU caches coherent at the
granularity of cache lines (Lesson~10),\index{cache coherence} units of a
few tens of bytes; this is affordable because the hardware coordinates the
caches within nanoseconds.  In a DSM every act of coordination costs
messages over the network, which makes such a fine granularity far too
expensive.  The candidates are:
\begin{description}
  \item[Variables] Sharing individual variables is too fine as well.  A
    variable occupies only a few bytes, so the messages needed to
    coordinate each access would cost far more than the data are worth.
  \item[Pages] A page, commonly \SI{4}{\kibi\byte}, is the unit that the
    operating system already manages.\margin{A \SI{2}{\mebi\byte} large page
    (Lesson~8) would cover 512 base pages at once: fewer and larger
    transfers, and correspondingly more false sharing.}  A page-based DSM
    is implemented by the operating system in cooperation with the MMU
    (\cref{sec:l15-implementing}); it is transparent to applications and
    independent of the language they are written in.
  \item[Objects] A language runtime knows the objects that a program
    creates.  An object-based DSM, implemented in the runtime, shares units
    that correspond to the program's data structures and needs no change
    to the operating system, but only programs that use that runtime
    benefit from it.
\end{description}

A coarser granularity amortizes the cost of a transfer or of an act of
coordination over more data, but it also places unrelated data in the same
unit.\caution{False sharing also occurs on a single machine, when two
variables share a cache line and the coherence protocol (Lesson~10) moves
the line back and forth.  A \SI{4}{\kibi\byte} page holds 64 lines of
\SI{64}{\byte}, and each move of it is a network transfer rather than a
cache miss.}

\begin{definition}[False sharing]
  The situation in which data that are not shared---each accessed by
  processes on a different node---reside in the same unit of sharing, so
  that the DSM coordinates the accesses to that unit as if the data were
  shared, is called \term{false sharing}[偽共有].
\end{definition}

\begin{figure}[htbp]
  \centering
  % False sharing: (a) x and y on the same page p, written by processes on
% different nodes, make p move on every write; (b) on separate pages p and
% q, nothing is transferred.
% Usage: % False sharing: (a) x and y on the same page p, written by processes on
% different nodes, make p move on every write; (b) on separate pages p and
% q, nothing is transferred.
% Usage: % False sharing: (a) x and y on the same page p, written by processes on
% different nodes, make p move on every write; (b) on separate pages p and
% q, nothing is transferred.
% Usage: \input{figures/l15-false-sharing}   (width ~118 mm)
\begin{tikzpicture}[x=1mm, y=1mm, font=\scriptsize\sffamily,
  >={Stealth[length=1.4mm]},
  mach/.style={draw, rounded corners=0.8mm, fill=white},
  half/.style={draw, fill=ln-accent!20, minimum width=6mm, minimum height=5mm, inner sep=0pt, font=\scriptsize\ttfamily},
  empty/.style={draw, fill=white, minimum width=6mm, minimum height=5mm, inner sep=0pt},
  gone/.style={draw, densely dashed, ln-muted, minimum width=12mm, minimum height=5mm, inner sep=0pt},
  lab/.style={font=\scriptsize\sffamily, inner sep=0.8pt},
  who/.style={font=\scriptsize\sffamily\bfseries, anchor=north west, inner sep=1pt},
  ttl/.style={font=\footnotesize\sffamily\bfseries, anchor=west}]
  % ---------------- (a) same page
  \node[ttl] at (0,23) {\iflangja{(a) \texttt{x} と \texttt{y} が同じページ}{(a) \texttt{x} and \texttt{y} on one page}};
  \draw[mach] (0,0) rectangle (21,17);
  \node[who] at (0.5,16.5) {\iflangja{ノード 1}{node 1}};
  \node[lab] at (10.5,11) {\iflangja{\texttt{x} に書く}{writes \texttt{x}}};
  \node[half] at (7.5,4.5) {x};
  \node[half] at (13.5,4.5) {y};
  \node[lab, anchor=west] at (16.8,4.5) {$p$};
  \draw[mach] (35,0) rectangle (56,17);
  \node[who] at (35.5,16.5) {\iflangja{ノード 2}{node 2}};
  \node[lab] at (45.5,11) {\iflangja{\texttt{y} に書く}{writes \texttt{y}}};
  \node[gone] at (45.5,4.5) {};
  \draw[->, thick, ln-caution] (21,13) to[bend left=35] node[lab, above, text=ln-caution] {$p$} (35,13);
  \draw[->, thick, ln-caution] (35,5) to[bend left=35] node[lab, below, text=ln-caution] {$p$} (21,5);
  % ---------------- (b) separate pages
  \begin{scope}[shift={(62,0)}]
    \node[ttl] at (0,23) {\iflangja{(b) 別々のページ}{(b) separate pages}};
    \draw[mach] (0,0) rectangle (21,17);
    \node[who] at (0.5,16.5) {\iflangja{ノード 1}{node 1}};
    \node[lab] at (10.5,11) {\iflangja{\texttt{x} に書く}{writes \texttt{x}}};
    \node[half] at (7.5,4.5) {x};
    \node[empty] at (13.5,4.5) {};
    \node[lab, anchor=west] at (16.8,4.5) {$p$};
    \draw[mach] (35,0) rectangle (56,17);
    \node[who] at (35.5,16.5) {\iflangja{ノード 2}{node 2}};
    \node[lab] at (45.5,11) {\iflangja{\texttt{y} に書く}{writes \texttt{y}}};
    \node[half] at (42.5,4.5) {y};
    \node[empty] at (48.5,4.5) {};
    \node[lab, anchor=west] at (51.8,4.5) {$q$};
    \node[lab, text=ln-muted, align=center] at (28,9) {\iflangja{転送\\なし}{no\\transfers}};
  \end{scope}
\end{tikzpicture}
   (width ~118 mm)
\begin{tikzpicture}[x=1mm, y=1mm, font=\scriptsize\sffamily,
  >={Stealth[length=1.4mm]},
  mach/.style={draw, rounded corners=0.8mm, fill=white},
  half/.style={draw, fill=ln-accent!20, minimum width=6mm, minimum height=5mm, inner sep=0pt, font=\scriptsize\ttfamily},
  empty/.style={draw, fill=white, minimum width=6mm, minimum height=5mm, inner sep=0pt},
  gone/.style={draw, densely dashed, ln-muted, minimum width=12mm, minimum height=5mm, inner sep=0pt},
  lab/.style={font=\scriptsize\sffamily, inner sep=0.8pt},
  who/.style={font=\scriptsize\sffamily\bfseries, anchor=north west, inner sep=1pt},
  ttl/.style={font=\footnotesize\sffamily\bfseries, anchor=west}]
  % ---------------- (a) same page
  \node[ttl] at (0,23) {\iflangja{(a) \texttt{x} と \texttt{y} が同じページ}{(a) \texttt{x} and \texttt{y} on one page}};
  \draw[mach] (0,0) rectangle (21,17);
  \node[who] at (0.5,16.5) {\iflangja{ノード 1}{node 1}};
  \node[lab] at (10.5,11) {\iflangja{\texttt{x} に書く}{writes \texttt{x}}};
  \node[half] at (7.5,4.5) {x};
  \node[half] at (13.5,4.5) {y};
  \node[lab, anchor=west] at (16.8,4.5) {$p$};
  \draw[mach] (35,0) rectangle (56,17);
  \node[who] at (35.5,16.5) {\iflangja{ノード 2}{node 2}};
  \node[lab] at (45.5,11) {\iflangja{\texttt{y} に書く}{writes \texttt{y}}};
  \node[gone] at (45.5,4.5) {};
  \draw[->, thick, ln-caution] (21,13) to[bend left=35] node[lab, above, text=ln-caution] {$p$} (35,13);
  \draw[->, thick, ln-caution] (35,5) to[bend left=35] node[lab, below, text=ln-caution] {$p$} (21,5);
  % ---------------- (b) separate pages
  \begin{scope}[shift={(62,0)}]
    \node[ttl] at (0,23) {\iflangja{(b) 別々のページ}{(b) separate pages}};
    \draw[mach] (0,0) rectangle (21,17);
    \node[who] at (0.5,16.5) {\iflangja{ノード 1}{node 1}};
    \node[lab] at (10.5,11) {\iflangja{\texttt{x} に書く}{writes \texttt{x}}};
    \node[half] at (7.5,4.5) {x};
    \node[empty] at (13.5,4.5) {};
    \node[lab, anchor=west] at (16.8,4.5) {$p$};
    \draw[mach] (35,0) rectangle (56,17);
    \node[who] at (35.5,16.5) {\iflangja{ノード 2}{node 2}};
    \node[lab] at (45.5,11) {\iflangja{\texttt{y} に書く}{writes \texttt{y}}};
    \node[half] at (42.5,4.5) {y};
    \node[empty] at (48.5,4.5) {};
    \node[lab, anchor=west] at (51.8,4.5) {$q$};
    \node[lab, text=ln-muted, align=center] at (28,9) {\iflangja{転送\\なし}{no\\transfers}};
  \end{scope}
\end{tikzpicture}
   (width ~118 mm)
\begin{tikzpicture}[x=1mm, y=1mm, font=\scriptsize\sffamily,
  >={Stealth[length=1.4mm]},
  mach/.style={draw, rounded corners=0.8mm, fill=white},
  half/.style={draw, fill=ln-accent!20, minimum width=6mm, minimum height=5mm, inner sep=0pt, font=\scriptsize\ttfamily},
  empty/.style={draw, fill=white, minimum width=6mm, minimum height=5mm, inner sep=0pt},
  gone/.style={draw, densely dashed, ln-muted, minimum width=12mm, minimum height=5mm, inner sep=0pt},
  lab/.style={font=\scriptsize\sffamily, inner sep=0.8pt},
  who/.style={font=\scriptsize\sffamily\bfseries, anchor=north west, inner sep=1pt},
  ttl/.style={font=\footnotesize\sffamily\bfseries, anchor=west}]
  % ---------------- (a) same page
  \node[ttl] at (0,23) {\iflangja{(a) \texttt{x} と \texttt{y} が同じページ}{(a) \texttt{x} and \texttt{y} on one page}};
  \draw[mach] (0,0) rectangle (21,17);
  \node[who] at (0.5,16.5) {\iflangja{ノード 1}{node 1}};
  \node[lab] at (10.5,11) {\iflangja{\texttt{x} に書く}{writes \texttt{x}}};
  \node[half] at (7.5,4.5) {x};
  \node[half] at (13.5,4.5) {y};
  \node[lab, anchor=west] at (16.8,4.5) {$p$};
  \draw[mach] (35,0) rectangle (56,17);
  \node[who] at (35.5,16.5) {\iflangja{ノード 2}{node 2}};
  \node[lab] at (45.5,11) {\iflangja{\texttt{y} に書く}{writes \texttt{y}}};
  \node[gone] at (45.5,4.5) {};
  \draw[->, thick, ln-caution] (21,13) to[bend left=35] node[lab, above, text=ln-caution] {$p$} (35,13);
  \draw[->, thick, ln-caution] (35,5) to[bend left=35] node[lab, below, text=ln-caution] {$p$} (21,5);
  % ---------------- (b) separate pages
  \begin{scope}[shift={(62,0)}]
    \node[ttl] at (0,23) {\iflangja{(b) 別々のページ}{(b) separate pages}};
    \draw[mach] (0,0) rectangle (21,17);
    \node[who] at (0.5,16.5) {\iflangja{ノード 1}{node 1}};
    \node[lab] at (10.5,11) {\iflangja{\texttt{x} に書く}{writes \texttt{x}}};
    \node[half] at (7.5,4.5) {x};
    \node[empty] at (13.5,4.5) {};
    \node[lab, anchor=west] at (16.8,4.5) {$p$};
    \draw[mach] (35,0) rectangle (56,17);
    \node[who] at (35.5,16.5) {\iflangja{ノード 2}{node 2}};
    \node[lab] at (45.5,11) {\iflangja{\texttt{y} に書く}{writes \texttt{y}}};
    \node[half] at (42.5,4.5) {y};
    \node[empty] at (48.5,4.5) {};
    \node[lab, anchor=west] at (51.8,4.5) {$q$};
    \node[lab, text=ln-muted, align=center] at (28,9) {\iflangja{転送\\なし}{no\\transfers}};
  \end{scope}
\end{tikzpicture}

  \caption{False sharing.  (a) Node 1 writes \texttt{x} and node 2 writes
    \texttt{y}; since both reside on page $p$, the page moves back and
    forth.  (b) On separate pages $p$ and $q$, both accesses remain local.}
  \label{fig:l15-false-sharing}
\end{figure}

\begin{example}
  Variables \texttt{x} and \texttt{y} reside on the same page $p$.  A
  process on node~1 writes only \texttt{x}, a process on node~2 writes
  only \texttt{y}; each writes its variable every \SI{10}{\milli\second},
  and the writes of the two processes alternate.  The processes share no
  data, yet a page-based DSM sees two nodes writing page $p$
  (\cref{fig:l15-false-sharing}a).  Every write finds the current contents
  of $p$ on the other node, so $p$ is transferred every
  \SI{5}{\milli\second}, 200 times per second, which amounts to
  $200\times\SI{4}{\kibi\byte} = \SI{800}{\kibi\byte}$ of network traffic
  per second.  If a transfer takes \SI{0.5}{\milli\second}, each process
  waits \SI{0.5}{\milli\second} of every \SI{10}{\milli\second},
  \SI{5}{\percent} of its time.  With \texttt{x} and \texttt{y} on
  different pages (\cref{fig:l15-false-sharing}b), no page is transferred
  after the first access.
\end{example}

False sharing is avoided by placing data that different nodes access in
different units, which the programmer, the compiler or a memory allocator
can do when the access pattern is known.

\section{Access algorithms}
\label{sec:l15-access}

What a DSM must provide depends on the access algorithm of the applications
that run on it.\source{P4L3, DSM design: access algorithm; Nitzberg and Lo.}
The access algorithm describes how many nodes read and write the same unit
of memory at the same time; three algorithms are distinguished.
\begin{description}
  \item[SRSW] Under
    \term{single reader/single writer}[単一リーダ単一ライタ]%
    \index{SRSW|see{single reader/single writer}} (SRSW) access, only one
    node at a time reads and writes a unit.  The main task of the DSM is to
    make memory that resides on other nodes accessible; since no two nodes
    use a unit at the same time, there are no concurrent accesses to
    coordinate.
  \item[MRSW] Under
    \term{multiple readers/single writer}[複数リーダ単一ライタ]%
    \index{MRSW|see{multiple readers/single writer}} (MRSW) access, several
    nodes may read a unit at the same time, but only one node writes it.
    The DSM must ensure that every read returns the value of the most
    recent write.
  \item[MRMW] Under
    \term{multiple readers/multiple writers}[複数リーダ複数ライタ]%
    \index{MRMW|see{multiple readers/multiple writers}} (MRMW) access,
    several nodes may both read and write a unit.  In addition, the DSM
    must order the writes correctly, so that all nodes observe them in the
    same order and reads return the most recent write.\caution{The access
    algorithm states what the applications do, not a discipline that the
    DSM enforces: under MRMW the DSM must cope with concurrent writers, not
    prevent them.}
\end{description}

\section{Migration and replication}
\label{sec:l15-migration}

The key performance metric of a DSM is access latency, the time that a
read or write of shared memory takes.\source{P4L3, DSM design: migration
vs.\ replication; Nitzberg and Lo.}  Since a remote access is orders of
magnitude slower than a local one, even a small fraction of remote accesses
raises the average sharply: one remote access in ten thousand already adds
half as much again, and one in a hundred multiplies the average by about
fifty.

\begin{example}
  \label{ex:l15-latency}
  A local memory access takes \SI{100}{\nano\second}, and an access to a
  page that must first be obtained from another node takes
  \SI{0.5}{\milli\second}.  If a fraction $f$ of the accesses are remote,
  the average access latency is
  \begin{equation}
    \bar{t} = (1-f)\times\SI{100}{\nano\second} + f\times\SI{0.5}{\milli\second} .
    \label{eq:l15-latency}
  \end{equation}
  For $f=10^{-4}$, $\bar{t} \approx \SI{100}{\nano\second} +
  \SI{50}{\nano\second} = \SI{150}{\nano\second}$, 1.5 times the latency of
  local memory; for $f=10^{-2}$, $\bar{t} = \SI{99}{\nano\second} +
  \SI{5000}{\nano\second} \approx \SI{5.1}{\micro\second}$, about 51 times.
\end{example}

Low latency therefore requires locality: the memory that a process accesses
should reside on the process's own node.\margin{The
\SI{0.5}{\milli\second} of \cref{ex:l15-latency} is a conservative figure.
A \SI{4}{\kibi\byte} page is on the wire for \SI{0.33}{\micro\second} at
\SI{100}{\giga\bit\per\second}; with the fault, the request and the
protocol overhead, a fetch over such an interconnect costs a few
microseconds.}  Two techniques bring memory to
the processes that access it (\cref{fig:l15-migration}).
\begin{description}
  \item[Migration] The DSM keeps a single copy of each unit and moves it to
    the node that accesses it.  With one copy there is nothing to keep
    consistent.  Migration suits SRSW access, where one node uses a unit
    for a stretch of accesses before another node needs it.  It transfers
    the unit whenever a different node accesses it, which is expensive when
    nodes take turns for short periods, and it cannot serve several nodes at
    the same time.
  \item[Replication] The DSM copies a unit to several nodes,\tip{Count
    messages: with $n$ replicas a write costs $n-1$ invalidations, and the
    next access on each invalidated node costs a fetch.  Replication pays
    while the local reads it saves outnumber them.} each of which
    accesses its local replica; caching is a form of
    replication\index{replication}.
    Replication is the more general technique: it lets several nodes read,
    or even write, a unit at the same time, as MRSW and MRMW access
    require.  Its price is consistency management
    (\cref{sec:l15-management}): when one replica is written, the others
    must be invalidated or updated.  This overhead grows with the number of
    replicas, and the DSM can control it by limiting the number of replicas
    of each unit.
\end{description}

\begin{figure}[htbp]
  \centering
  % Migration and replication: (a) the single copy of page p moves to node B,
% which accesses it; (b) nodes B and C hold replicas of p, and a write by A
% makes the DSM invalidate them.
% Usage: % Migration and replication: (a) the single copy of page p moves to node B,
% which accesses it; (b) nodes B and C hold replicas of p, and a write by A
% makes the DSM invalidate them.
% Usage: % Migration and replication: (a) the single copy of page p moves to node B,
% which accesses it; (b) nodes B and C hold replicas of p, and a write by A
% makes the DSM invalidate them.
% Usage: \input{figures/l15-migration}   (width ~118 mm)
\begin{tikzpicture}[x=1mm, y=1mm, font=\scriptsize\sffamily,
  >={Stealth[length=1.4mm]},
  mach/.style={draw, rounded corners=0.8mm, fill=white},
  own/.style={draw, fill=ln-accent!40, minimum width=9mm, minimum height=5mm, inner sep=0pt},
  rep/.style={draw, fill=ln-accent!12, minimum width=9mm, minimum height=5mm, inner sep=0pt},
  gone/.style={draw, densely dashed, ln-muted, minimum width=9mm, minimum height=5mm, inner sep=0pt},
  lab/.style={font=\scriptsize\sffamily, inner sep=0.8pt},
  who/.style={font=\scriptsize\sffamily\bfseries, anchor=north west, inner sep=1pt},
  ttl/.style={font=\footnotesize\sffamily\bfseries, anchor=west}]
  % \lnnodes: three node boxes A, B, C at x = 0, 20, 40
  \def\lnnodes{%
    \foreach \nx/\nn in {0/A,20/B,40/C} {
      \draw[mach] (\nx,0) rectangle ++(16,13);
      \node[who] at (\nx+0.4,12.6) {\nn};
    }}
  % ---------------- (a) migration
  \node[ttl] at (0,26) {\iflangja{(a) マイグレーション}{(a) migration}};
  \lnnodes
  \node[gone] at (8,5) {};
  \node[own] at (28,5) {$p$};
  \draw[->, thick, ln-accent] (8,13) to[bend left=40] node[lab, above, text=ln-accent] {\iflangja{$p$ を移動}{move $p$}} (28,13);
  \node[lab, anchor=north] at (28,-1) {\iflangja{$p$ にアクセス}{accesses $p$}};
  % ---------------- (b) replication
  \begin{scope}[shift={(62,0)}]
    \node[ttl] at (0,26) {\iflangja{(b) レプリケーション}{(b) replication}};
    \lnnodes
    \node[own] at (8,5) {$p$};
    \node[rep] at (28,5) {$p$};
    \node[rep] at (48,5) {$p$};
    \draw[->, thick, ln-caution] (10,13) to[bend left=40] (28,13);
    \draw[->, thick, ln-caution] (6,13) to[bend left=38] node[lab, above=0.8mm, fill=white, text=ln-caution] {\iflangja{無効化}{invalidate}} (48,13);
    \node[lab, anchor=north] at (8,-1) {\iflangja{$p$ に書く}{writes $p$}};
    \node[lab, anchor=north] at (38,-1) {\iflangja{$p$ を読む}{read $p$}};
  \end{scope}
\end{tikzpicture}
   (width ~118 mm)
\begin{tikzpicture}[x=1mm, y=1mm, font=\scriptsize\sffamily,
  >={Stealth[length=1.4mm]},
  mach/.style={draw, rounded corners=0.8mm, fill=white},
  own/.style={draw, fill=ln-accent!40, minimum width=9mm, minimum height=5mm, inner sep=0pt},
  rep/.style={draw, fill=ln-accent!12, minimum width=9mm, minimum height=5mm, inner sep=0pt},
  gone/.style={draw, densely dashed, ln-muted, minimum width=9mm, minimum height=5mm, inner sep=0pt},
  lab/.style={font=\scriptsize\sffamily, inner sep=0.8pt},
  who/.style={font=\scriptsize\sffamily\bfseries, anchor=north west, inner sep=1pt},
  ttl/.style={font=\footnotesize\sffamily\bfseries, anchor=west}]
  % \lnnodes: three node boxes A, B, C at x = 0, 20, 40
  \def\lnnodes{%
    \foreach \nx/\nn in {0/A,20/B,40/C} {
      \draw[mach] (\nx,0) rectangle ++(16,13);
      \node[who] at (\nx+0.4,12.6) {\nn};
    }}
  % ---------------- (a) migration
  \node[ttl] at (0,26) {\iflangja{(a) マイグレーション}{(a) migration}};
  \lnnodes
  \node[gone] at (8,5) {};
  \node[own] at (28,5) {$p$};
  \draw[->, thick, ln-accent] (8,13) to[bend left=40] node[lab, above, text=ln-accent] {\iflangja{$p$ を移動}{move $p$}} (28,13);
  \node[lab, anchor=north] at (28,-1) {\iflangja{$p$ にアクセス}{accesses $p$}};
  % ---------------- (b) replication
  \begin{scope}[shift={(62,0)}]
    \node[ttl] at (0,26) {\iflangja{(b) レプリケーション}{(b) replication}};
    \lnnodes
    \node[own] at (8,5) {$p$};
    \node[rep] at (28,5) {$p$};
    \node[rep] at (48,5) {$p$};
    \draw[->, thick, ln-caution] (10,13) to[bend left=40] (28,13);
    \draw[->, thick, ln-caution] (6,13) to[bend left=38] node[lab, above=0.8mm, fill=white, text=ln-caution] {\iflangja{無効化}{invalidate}} (48,13);
    \node[lab, anchor=north] at (8,-1) {\iflangja{$p$ に書く}{writes $p$}};
    \node[lab, anchor=north] at (38,-1) {\iflangja{$p$ を読む}{read $p$}};
  \end{scope}
\end{tikzpicture}
   (width ~118 mm)
\begin{tikzpicture}[x=1mm, y=1mm, font=\scriptsize\sffamily,
  >={Stealth[length=1.4mm]},
  mach/.style={draw, rounded corners=0.8mm, fill=white},
  own/.style={draw, fill=ln-accent!40, minimum width=9mm, minimum height=5mm, inner sep=0pt},
  rep/.style={draw, fill=ln-accent!12, minimum width=9mm, minimum height=5mm, inner sep=0pt},
  gone/.style={draw, densely dashed, ln-muted, minimum width=9mm, minimum height=5mm, inner sep=0pt},
  lab/.style={font=\scriptsize\sffamily, inner sep=0.8pt},
  who/.style={font=\scriptsize\sffamily\bfseries, anchor=north west, inner sep=1pt},
  ttl/.style={font=\footnotesize\sffamily\bfseries, anchor=west}]
  % \lnnodes: three node boxes A, B, C at x = 0, 20, 40
  \def\lnnodes{%
    \foreach \nx/\nn in {0/A,20/B,40/C} {
      \draw[mach] (\nx,0) rectangle ++(16,13);
      \node[who] at (\nx+0.4,12.6) {\nn};
    }}
  % ---------------- (a) migration
  \node[ttl] at (0,26) {\iflangja{(a) マイグレーション}{(a) migration}};
  \lnnodes
  \node[gone] at (8,5) {};
  \node[own] at (28,5) {$p$};
  \draw[->, thick, ln-accent] (8,13) to[bend left=40] node[lab, above, text=ln-accent] {\iflangja{$p$ を移動}{move $p$}} (28,13);
  \node[lab, anchor=north] at (28,-1) {\iflangja{$p$ にアクセス}{accesses $p$}};
  % ---------------- (b) replication
  \begin{scope}[shift={(62,0)}]
    \node[ttl] at (0,26) {\iflangja{(b) レプリケーション}{(b) replication}};
    \lnnodes
    \node[own] at (8,5) {$p$};
    \node[rep] at (28,5) {$p$};
    \node[rep] at (48,5) {$p$};
    \draw[->, thick, ln-caution] (10,13) to[bend left=40] (28,13);
    \draw[->, thick, ln-caution] (6,13) to[bend left=38] node[lab, above=0.8mm, fill=white, text=ln-caution] {\iflangja{無効化}{invalidate}} (48,13);
    \node[lab, anchor=north] at (8,-1) {\iflangja{$p$ に書く}{writes $p$}};
    \node[lab, anchor=north] at (38,-1) {\iflangja{$p$ を読む}{read $p$}};
  \end{scope}
\end{tikzpicture}

  \caption{Bringing page $p$ to the nodes that access it.  (a) Migration
    moves the single copy.  (b) Replication keeps copies on several nodes,
    which read $p$ locally; a write by the owner A invalidates the replicas
    on B and C (dashed).}
  \label{fig:l15-migration}
\end{figure}

\begin{example}
  Page $p$ initially resides on node~A, and nodes A, B and C access it in
  the sequence of \cref{tab:l15-migration}.  A page transfer takes
  \SI{0.5}{\milli\second} and an invalidation message
  \SI{0.05}{\milli\second}.\margin{An invalidation carries only the
  identifier of a page, a few tens of bytes, so its cost is that of a
  message exchange and not of the data.}  Migration moves $p$ on every
  access by a node that does not hold it: 7 transfers,
  \SI{3.5}{\milli\second}.  With
  replication, B and C fetch replicas once and read them locally until A's
  write in step~6 invalidates them: 4 transfers and 2 invalidations,
  \SI{2.1}{\milli\second}.  Had A made all eight accesses, neither technique
  would transfer anything.  Had A and B instead alternated writes, every
  write would invalidate the other node's replica and every next access
  would fetch it again, so that replication would cost as many transfers as
  migration plus the invalidations.
\end{example}

\begin{table}[htbp]
  \centering
  \caption{Accesses to page $p$ under migration and under replication with
    invalidation on write.}
  \label{tab:l15-migration}
  \small
  \begin{tabular}{@{}rl ll ll@{}}
    \toprule
    & & \multicolumn{2}{c}{Migration} & \multicolumn{2}{c@{}}{Replication} \\
    \cmidrule(lr){3-4}\cmidrule(l){5-6}
    Step & Access & $p$ at & Transfer & Copies at & Messages \\
    \midrule
    0 & (initial) & A & --- & A & --- \\
    1 & A writes & A & --- & A & --- \\
    2 & B reads & B & A$\to$B & A, B & fetch from A \\
    3 & C reads & C & B$\to$C & A, B, C & fetch from A \\
    4 & B reads & B & C$\to$B & A, B, C & --- \\
    5 & C reads & C & B$\to$C & A, B, C & --- \\
    6 & A writes & A & C$\to$A & A & invalidate B, C \\
    7 & B reads & B & A$\to$B & A, B & fetch from A \\
    8 & C reads & C & B$\to$C & A, B, C & fetch from A \\
    \midrule
    \multicolumn{2}{@{}l}{Total} & & 7 transfers & & 4 fetches, 2 invalidations \\
    \bottomrule
  \end{tabular}
\end{table}

Neither technique is best for every workload.  Migration is appropriate
when a unit is used by one node at a time for long stretches.  Replication
is necessary when several nodes access a unit concurrently, and it pays off
when reads dominate.\margin{Replication with invalidation on write is the
write-invalidate strategy of Lesson~10 in software: a page instead of a
cache line, a message instead of a bus transaction.}  When replicated units
are written often, the consistency management that every write triggers can
cost more than the locality of the replicas saves.

\begin{keypoint}
  DSM performance is governed by access latency, and low latency requires
  locality.  Migration moves a single copy and suits SRSW access;
  replication keeps several copies, supports concurrent access, and pays
  for it with consistency management that grows with the number of
  replicas and the frequency of writes.
\end{keypoint}

\section{Consistency management}
\label{sec:l15-management}

Once memory is replicated, the DSM must keep the replicas
consistent.\source{P4L3, DSM design: consistency management; Nitzberg and
Lo.}  On a shared memory multiprocessor, software orders its accesses with
locks, and the hardware keeps the copies of memory in the CPU caches
coherent (Lesson~10).  A DSM has no such hardware; its mechanisms resemble
those with which a distributed file system keeps client caches consistent
(Lesson~14):
\begin{description}
  \item[Push] When a replica is written, the DSM immediately pushes
    invalidations to the nodes that hold other replicas of the same unit.
    This approach is \emph{proactive} and \emph{eager}; it is also
    \emph{pessimistic}, because it assumes that a stale replica would be
    read and cause harm, and removes stale replicas as soon as they
    arise.\caution{Consistency management in a DSM is not cache coherence:
    it is carried out by software, at the granularity of a page or an
    object, and only at the points the consistency model requires.}
  \item[Pull] Nodes pull information about modifications from the owner of
    a unit, periodically or when they need it.  This approach is
    \emph{reactive} (on demand) and \emph{lazy}; it is also
    \emph{optimistic}, because it hopes that using a stale replica for a
    while does no harm, and it saves the messages of eager propagation.
\end{description}
When the DSM actually triggers these mechanisms---on every write, on every
access, periodically, or only at points that the application
designates---is determined by the consistency model that the DSM provides
for the shared state (\cref{sec:l15-models}).\margin{Push and pull are the
DSM counterparts of the server-driven and client-driven cache consistency
of Lesson~14.}

\begin{figure}[htbp]
  \centering
  % Consistency management: node A writes page p at 0.5 s, node B holds a
% replica and reads p at 1 s and 3 s.  (a) Push: A invalidates B's replica
% at once.  (b) Pull: B polls A every 2 s, and its replica is stale in
% between.
% Usage: % Consistency management: node A writes page p at 0.5 s, node B holds a
% replica and reads p at 1 s and 3 s.  (a) Push: A invalidates B's replica
% at once.  (b) Pull: B polls A every 2 s, and its replica is stale in
% between.
% Usage: % Consistency management: node A writes page p at 0.5 s, node B holds a
% replica and reads p at 1 s and 3 s.  (a) Push: A invalidates B's replica
% at once.  (b) Pull: B polls A every 2 s, and its replica is stale in
% between.
% Usage: \input{figures/l15-push-pull}   (width ~116 mm)
\begin{tikzpicture}[x=1mm, y=1mm, font=\scriptsize\sffamily,
  >={Stealth[length=1.3mm]},
  row/.style={font=\scriptsize\sffamily, anchor=east},
  lab/.style={font=\scriptsize\sffamily, inner sep=0.8pt},
  ev/.style={circle, fill=black, inner sep=0pt, minimum size=1.3mm},
  ttl/.style={font=\footnotesize\sffamily\bfseries, anchor=west}]
  % time t (s) -> x = 22 + 23 t (mm); rows: A at 9, B at 0 (inside a scope)
  \def\lnrows{%
    \node[row] at (20,9) {\iflangja{A（書き手）}{A (writer)}};
    \node[row] at (20,0) {\iflangja{B（複製）}{B (replica)}};
    \foreach \y in {0,9} { \draw[ln-rule] (22,\y) -- (116,\y); }
    \node[ev] at ({22+23*0.5},9) {};
    \node[lab, anchor=south] at ({22+23*0.5},10) {\iflangja{書く}{write}};
    \foreach \t in {1,3} { \node[ev] at ({22+23*\t},0) {}; }
  }
  % ---------------- (a) push
  \begin{scope}[shift={(0,27)}]
    \node[ttl] at (0,16) {\iflangja{(a) プッシュ}{(a) push}};
    \lnrows
    \draw[->, thick, ln-caution] ({22+23*0.5},8) -- ({22+23*0.5},1);
    \node[lab, anchor=east, text=ln-caution] at ({21.4+23*0.5},4.5) {\iflangja{無効化}{invalidate}};
    \draw[->, densely dashed] ({22+23*1},8) -- ({22+23*1},1);
    \node[lab, anchor=west] at ({22.6+23*1},4.5) {\iflangja{取得}{fetch}};
    \node[lab, anchor=north] at ({22+23*1},-1) {\iflangja{読む：新}{read: new}};
    \node[lab, anchor=north] at ({22+23*3},-1) {\iflangja{読む：新}{read: new}};
  \end{scope}
  % ---------------- (b) pull
  \begin{scope}[shift={(0,0)}]
    \node[ttl] at (0,16) {\iflangja{(b) プル（2 秒ごと）}{(b) pull (every 2 s)}};
    \lnrows
    \fill[ln-caution!22] ({22+23*0.5},-1) rectangle ({22+23*2},1);
    \draw[ln-rule] ({22+23*0.5},0) -- ({22+23*2},0);
    \foreach \t in {1,3} { \node[ev] at ({22+23*\t},0) {}; }
    \node[lab, anchor=south, text=ln-caution] at ({22+23*1.5},1.2) {\iflangja{古い}{stale}};
    \foreach \t in {0,2,4} {
      \draw[<->, densely dashed] ({22+23*\t},1) -- ({22+23*\t},8);
    }
    \node[lab, anchor=west] at ({22.6+23*2},4.5) {\iflangja{ポーリング}{poll}};
    \node[lab, anchor=north] at ({22+23*1},-1) {\iflangja{読む：古}{read: stale}};
    \node[lab, anchor=north] at ({22+23*3},-1) {\iflangja{読む：新}{read: new}};
  \end{scope}
  % time axis
  \draw[->] (22,-8) -- (117,-8);
  \foreach \t in {0,1,2,3,4} {
    \draw ({22+23*\t},-7.2) -- ({22+23*\t},-8.8);
    \node[lab, anchor=north] at ({22+23*\t},-9) {\t};
  }
  \node[lab, anchor=north east] at (20,-8.6) {\iflangja{時間（秒）}{time (s)}};
\end{tikzpicture}
   (width ~116 mm)
\begin{tikzpicture}[x=1mm, y=1mm, font=\scriptsize\sffamily,
  >={Stealth[length=1.3mm]},
  row/.style={font=\scriptsize\sffamily, anchor=east},
  lab/.style={font=\scriptsize\sffamily, inner sep=0.8pt},
  ev/.style={circle, fill=black, inner sep=0pt, minimum size=1.3mm},
  ttl/.style={font=\footnotesize\sffamily\bfseries, anchor=west}]
  % time t (s) -> x = 22 + 23 t (mm); rows: A at 9, B at 0 (inside a scope)
  \def\lnrows{%
    \node[row] at (20,9) {\iflangja{A（書き手）}{A (writer)}};
    \node[row] at (20,0) {\iflangja{B（複製）}{B (replica)}};
    \foreach \y in {0,9} { \draw[ln-rule] (22,\y) -- (116,\y); }
    \node[ev] at ({22+23*0.5},9) {};
    \node[lab, anchor=south] at ({22+23*0.5},10) {\iflangja{書く}{write}};
    \foreach \t in {1,3} { \node[ev] at ({22+23*\t},0) {}; }
  }
  % ---------------- (a) push
  \begin{scope}[shift={(0,27)}]
    \node[ttl] at (0,16) {\iflangja{(a) プッシュ}{(a) push}};
    \lnrows
    \draw[->, thick, ln-caution] ({22+23*0.5},8) -- ({22+23*0.5},1);
    \node[lab, anchor=east, text=ln-caution] at ({21.4+23*0.5},4.5) {\iflangja{無効化}{invalidate}};
    \draw[->, densely dashed] ({22+23*1},8) -- ({22+23*1},1);
    \node[lab, anchor=west] at ({22.6+23*1},4.5) {\iflangja{取得}{fetch}};
    \node[lab, anchor=north] at ({22+23*1},-1) {\iflangja{読む：新}{read: new}};
    \node[lab, anchor=north] at ({22+23*3},-1) {\iflangja{読む：新}{read: new}};
  \end{scope}
  % ---------------- (b) pull
  \begin{scope}[shift={(0,0)}]
    \node[ttl] at (0,16) {\iflangja{(b) プル（2 秒ごと）}{(b) pull (every 2 s)}};
    \lnrows
    \fill[ln-caution!22] ({22+23*0.5},-1) rectangle ({22+23*2},1);
    \draw[ln-rule] ({22+23*0.5},0) -- ({22+23*2},0);
    \foreach \t in {1,3} { \node[ev] at ({22+23*\t},0) {}; }
    \node[lab, anchor=south, text=ln-caution] at ({22+23*1.5},1.2) {\iflangja{古い}{stale}};
    \foreach \t in {0,2,4} {
      \draw[<->, densely dashed] ({22+23*\t},1) -- ({22+23*\t},8);
    }
    \node[lab, anchor=west] at ({22.6+23*2},4.5) {\iflangja{ポーリング}{poll}};
    \node[lab, anchor=north] at ({22+23*1},-1) {\iflangja{読む：古}{read: stale}};
    \node[lab, anchor=north] at ({22+23*3},-1) {\iflangja{読む：新}{read: new}};
  \end{scope}
  % time axis
  \draw[->] (22,-8) -- (117,-8);
  \foreach \t in {0,1,2,3,4} {
    \draw ({22+23*\t},-7.2) -- ({22+23*\t},-8.8);
    \node[lab, anchor=north] at ({22+23*\t},-9) {\t};
  }
  \node[lab, anchor=north east] at (20,-8.6) {\iflangja{時間（秒）}{time (s)}};
\end{tikzpicture}
   (width ~116 mm)
\begin{tikzpicture}[x=1mm, y=1mm, font=\scriptsize\sffamily,
  >={Stealth[length=1.3mm]},
  row/.style={font=\scriptsize\sffamily, anchor=east},
  lab/.style={font=\scriptsize\sffamily, inner sep=0.8pt},
  ev/.style={circle, fill=black, inner sep=0pt, minimum size=1.3mm},
  ttl/.style={font=\footnotesize\sffamily\bfseries, anchor=west}]
  % time t (s) -> x = 22 + 23 t (mm); rows: A at 9, B at 0 (inside a scope)
  \def\lnrows{%
    \node[row] at (20,9) {\iflangja{A（書き手）}{A (writer)}};
    \node[row] at (20,0) {\iflangja{B（複製）}{B (replica)}};
    \foreach \y in {0,9} { \draw[ln-rule] (22,\y) -- (116,\y); }
    \node[ev] at ({22+23*0.5},9) {};
    \node[lab, anchor=south] at ({22+23*0.5},10) {\iflangja{書く}{write}};
    \foreach \t in {1,3} { \node[ev] at ({22+23*\t},0) {}; }
  }
  % ---------------- (a) push
  \begin{scope}[shift={(0,27)}]
    \node[ttl] at (0,16) {\iflangja{(a) プッシュ}{(a) push}};
    \lnrows
    \draw[->, thick, ln-caution] ({22+23*0.5},8) -- ({22+23*0.5},1);
    \node[lab, anchor=east, text=ln-caution] at ({21.4+23*0.5},4.5) {\iflangja{無効化}{invalidate}};
    \draw[->, densely dashed] ({22+23*1},8) -- ({22+23*1},1);
    \node[lab, anchor=west] at ({22.6+23*1},4.5) {\iflangja{取得}{fetch}};
    \node[lab, anchor=north] at ({22+23*1},-1) {\iflangja{読む：新}{read: new}};
    \node[lab, anchor=north] at ({22+23*3},-1) {\iflangja{読む：新}{read: new}};
  \end{scope}
  % ---------------- (b) pull
  \begin{scope}[shift={(0,0)}]
    \node[ttl] at (0,16) {\iflangja{(b) プル（2 秒ごと）}{(b) pull (every 2 s)}};
    \lnrows
    \fill[ln-caution!22] ({22+23*0.5},-1) rectangle ({22+23*2},1);
    \draw[ln-rule] ({22+23*0.5},0) -- ({22+23*2},0);
    \foreach \t in {1,3} { \node[ev] at ({22+23*\t},0) {}; }
    \node[lab, anchor=south, text=ln-caution] at ({22+23*1.5},1.2) {\iflangja{古い}{stale}};
    \foreach \t in {0,2,4} {
      \draw[<->, densely dashed] ({22+23*\t},1) -- ({22+23*\t},8);
    }
    \node[lab, anchor=west] at ({22.6+23*2},4.5) {\iflangja{ポーリング}{poll}};
    \node[lab, anchor=north] at ({22+23*1},-1) {\iflangja{読む：古}{read: stale}};
    \node[lab, anchor=north] at ({22+23*3},-1) {\iflangja{読む：新}{read: new}};
  \end{scope}
  % time axis
  \draw[->] (22,-8) -- (117,-8);
  \foreach \t in {0,1,2,3,4} {
    \draw ({22+23*\t},-7.2) -- ({22+23*\t},-8.8);
    \node[lab, anchor=north] at ({22+23*\t},-9) {\t};
  }
  \node[lab, anchor=north east] at (20,-8.6) {\iflangja{時間（秒）}{time (s)}};
\end{tikzpicture}

  \caption{Node A writes page $p$ at \SI{0.5}{\second}; node B holds a
    replica.  (a) A pushes an invalidation, and B fetches the new contents
    on its next read.  (b) B polls every \SI{2}{\second}, and its replica is
    stale until the next poll.}
  \label{fig:l15-push-pull}
\end{figure}

\begin{example}
  Node~A owns page $p$, and node~B holds a replica.  A writes $p$ at
  $t=\SI{0.5}{\second}$, and B reads $p$ at \SI{1}{\second} and at
  \SI{3}{\second} (\cref{fig:l15-push-pull}).  With push, A sends B an
  invalidation at \SI{0.5}{\second}; B's read at \SI{1}{\second} finds no
  valid replica and fetches the new contents.  With pull, B asks A every
  \SI{2}{\second} whether $p$ has changed.  B's read at \SI{1}{\second}
  returns the stale contents; the poll at \SI{2}{\second} reveals the
  modification, and the read at \SI{3}{\second} returns the new contents.
  A replica may thus remain stale for up to one polling period.  In
  exchange, pulling costs messages at a fixed rate: 10 replicas that poll
  every \SI{2}{\second} send 5 messages per second, whether $p$ is written
  a hundred times per second or not at all.  Pushing sends no message while
  $p$ is not written, and at most one invalidation per replica for each
  write.\tip{Push costs messages in proportion to the writes, pull in
  proportion to time: compare the rate of writes with the rate of polling
  to see which is cheaper.}
\end{example}

\section{DSM architecture}
\label{sec:l15-architecture}

We now consider a page-based DSM implemented by the operating
system.\source{P4L3, DSM architecture, indexing distributed state;
Nitzberg and Lo.}  Its organization follows from the
preceding design choices.
\begin{itemize}
  \item Every node contributes part of its main memory, a number of page
    frames, to the DSM.  The contributed pages form a system-wide pool of
    shared memory, accessible from all nodes.
  \item Every node keeps local copies of pages that its processes access, a
    cache of replicas, to lower the access latency.
  \item Management is distributed as well: every node is responsible for
    part of the shared memory.
\end{itemize}
For each page the DSM distinguishes the following roles.
\begin{description}
  \item[Home node] The node that manages the page, also called its
    \emph{manager}.  It maintains the state of the page: which nodes access
    it, whether it has been modified, whether caching of it is enabled,
    whether it is locked, and where its replicas reside.  The address of a
    page names the node that manages it initially; management can be handed
    to another node without changing the address
    (\cref{sec:l15-indexing}).\margin{The home node plays the part that the
    server plays in a distributed file system (Lesson~14), except that
    every node is the home node of a portion of the memory.}
  \item[Owner] The node that currently holds the authoritative copy of the
    page and controls its updates, for example the node that most recently
    obtained the right to write it.  Ownership changes as the page is
    accessed, so the owner need not be the home node; the home node keeps
    track of the current owner.
  \item[Explicit replicas] Copies that the DSM creates deliberately, rather
    than as a side effect of access, for load balancing, for performance,
    or for reliability; in datacenters, for instance, data are commonly
    kept in three copies.\margin{Three is the default replication factor of
    HDFS, which places two replicas in one rack and the third in another
    rack (Hadoop documentation).}  The home node controls where these
    replicas are placed and how they are managed.
\end{description}

\subsection{Indexing distributed state}
\label{sec:l15-indexing}

To serve an access, a node must find the metadata of the page, which is
maintained by the page's manager.  A DSM organizes this metadata as
follows (\cref{fig:l15-indexing}).
\begin{itemize}
  \item Every page is identified by an address composed of a node ID and a
    physical frame number\index{physical frame number} (PFN, Lesson~8): the
    node ID names the node that contributes the frame and manages the page
    initially, and the PFN the frame on that node.
  \item A global map from page ID to the ID of the manager node is
    replicated on every node, so that any node finds the manager of any
    page locally, without exchanging messages.
  \item The metadata of the pages themselves are partitioned: each manager
    keeps the metadata of the pages it manages.\margin{Replicating the map
    is affordable because it is small: with the 8-bit node ID of the
    example below it has 256 entries.  The per-page metadata, far more
    numerous, stay partitioned.}
\end{itemize}
Rather than deriving the manager directly from bits of the page ID, the
global map can be implemented as a \emph{mapping table}: part of the page
ID serves as an index into the table, and the entry names the manager
node.  The manager of a set of pages can then be changed, for load
balancing or after a failure, by updating a table entry on all nodes, while
the page IDs, and every reference that uses them, remain unchanged.  Once
the table is interposed, the node-ID field is no more than an index into
it: the manager of a page is the node that the entry names, which need no
longer be the node the field denotes.\tip{The mapping table is to the
manager of a page what a page table is to a frame: a level of indirection
that lets what a name refers to change while the name stays the same.}

\begin{figure}[htbp]
  \centering
  % Indexing distributed state: the node-ID field of a page ID indexes the
% global mapping table, replicated on every node, which names the manager;
% the manager keeps the metadata of its pages.
% Usage: % Indexing distributed state: the node-ID field of a page ID indexes the
% global mapping table, replicated on every node, which names the manager;
% the manager keeps the metadata of its pages.
% Usage: % Indexing distributed state: the node-ID field of a page ID indexes the
% global mapping table, replicated on every node, which names the manager;
% the manager keeps the metadata of its pages.
% Usage: \input{figures/l15-indexing}   (width ~118 mm)
\begin{tikzpicture}[x=1mm, y=1mm, font=\scriptsize\sffamily,
  >={Stealth[length=1.4mm]},
  fld/.style={draw, minimum height=6mm, inner sep=0pt, font=\scriptsize\ttfamily},
  cellc/.style={draw, minimum height=5mm, inner sep=0pt},
  mach/.style={draw, rounded corners=0.8mm, fill=white},
  lab/.style={font=\scriptsize\sffamily, inner sep=0.8pt},
  hd/.style={font=\scriptsize\sffamily\bfseries, inner sep=0.8pt}]
  % page ID
  \node[hd, anchor=south west] at (0,37) {\iflangja{ページ ID}{page ID}};
  \node[fld, fill=ln-accent!20, minimum width=12mm] (nid) at (6,33) {05};
  \node[fld, fill=white, minimum width=22mm] at (23,33) {0012A3};
  \node[lab, anchor=north, text=ln-muted] at (6,29.6) {\iflangja{ノード ID}{node ID}};
  \node[lab, anchor=north, text=ln-muted] at (23,29.6) {PFN};
  % global mapping table
  \node[hd, anchor=south] at (58,37) {\iflangja{グローバルマップ}{global mapping table}};
  \node[lab, anchor=south, text=ln-muted] at (58,34) {\iflangja{（全ノードに複製）}{(replicated on every node)}};
  \node[lab, text=ln-muted] at (51,30.5) {\iflangja{索引}{index}};
  \node[lab, text=ln-muted] at (63,30.5) {\iflangja{管理ノード}{manager}};
  \foreach \ry/\ri/\rm in {26/04/4, 21/05/9, 16/06/6} {
    \node[cellc, minimum width=10mm, fill=white] at (51,\ry) {\ttfamily\ri};
    \node[cellc, minimum width=14mm, fill=white] at (63,\ry) {\rm};
  }
  \node[cellc, minimum width=14mm, fill=ln-accent!20] at (63,21) {9};
  \node[lab, text=ln-muted] at (51,11) {$\vdots$};
  \node[lab, text=ln-muted] at (63,11) {$\vdots$};
  \draw[->] (nid.south) ++(0,-3.6) -- (6,21) -- (46,21);
  % manager node
  \draw[mach] (78,4) rectangle (118,37);
  \node[hd, anchor=north west] at (78.5,36.5) {\iflangja{ノード 9（管理ノード）}{node 9 (manager)}};
  \node[lab, anchor=north west, text=ln-muted] at (78.5,32) {\iflangja{管理するページのメタデータ}{metadata of its pages}};
  \foreach \cx/\cw/\ct in {86.5/13/{\iflangja{ページ}{page}}, 99/12/{\iflangja{所有者}{owner}}, 110.5/11/{\iflangja{複製}{replicas}}} {
    \node[cellc, minimum width=\cw mm, fill=ln-box] at (\cx,24) {\ct};
  }
  \foreach \ry/\pa/\pb/\pc in {19/\ttfamily 050012A3/2/{4, 7}, 14/\ttfamily 050012A4/9/---} {
    \node[cellc, minimum width=13mm, fill=white] at (86.5,\ry) {\pa};
    \node[cellc, minimum width=12mm, fill=white] at (99,\ry) {\pb};
    \node[cellc, minimum width=11mm, fill=white] at (110.5,\ry) {\pc};
  }
  \node[lab, text=ln-muted] at (98,8.5) {$\vdots$};
  \draw[->] (70,21) -- (78,21);
\end{tikzpicture}
   (width ~118 mm)
\begin{tikzpicture}[x=1mm, y=1mm, font=\scriptsize\sffamily,
  >={Stealth[length=1.4mm]},
  fld/.style={draw, minimum height=6mm, inner sep=0pt, font=\scriptsize\ttfamily},
  cellc/.style={draw, minimum height=5mm, inner sep=0pt},
  mach/.style={draw, rounded corners=0.8mm, fill=white},
  lab/.style={font=\scriptsize\sffamily, inner sep=0.8pt},
  hd/.style={font=\scriptsize\sffamily\bfseries, inner sep=0.8pt}]
  % page ID
  \node[hd, anchor=south west] at (0,37) {\iflangja{ページ ID}{page ID}};
  \node[fld, fill=ln-accent!20, minimum width=12mm] (nid) at (6,33) {05};
  \node[fld, fill=white, minimum width=22mm] at (23,33) {0012A3};
  \node[lab, anchor=north, text=ln-muted] at (6,29.6) {\iflangja{ノード ID}{node ID}};
  \node[lab, anchor=north, text=ln-muted] at (23,29.6) {PFN};
  % global mapping table
  \node[hd, anchor=south] at (58,37) {\iflangja{グローバルマップ}{global mapping table}};
  \node[lab, anchor=south, text=ln-muted] at (58,34) {\iflangja{（全ノードに複製）}{(replicated on every node)}};
  \node[lab, text=ln-muted] at (51,30.5) {\iflangja{索引}{index}};
  \node[lab, text=ln-muted] at (63,30.5) {\iflangja{管理ノード}{manager}};
  \foreach \ry/\ri/\rm in {26/04/4, 21/05/9, 16/06/6} {
    \node[cellc, minimum width=10mm, fill=white] at (51,\ry) {\ttfamily\ri};
    \node[cellc, minimum width=14mm, fill=white] at (63,\ry) {\rm};
  }
  \node[cellc, minimum width=14mm, fill=ln-accent!20] at (63,21) {9};
  \node[lab, text=ln-muted] at (51,11) {$\vdots$};
  \node[lab, text=ln-muted] at (63,11) {$\vdots$};
  \draw[->] (nid.south) ++(0,-3.6) -- (6,21) -- (46,21);
  % manager node
  \draw[mach] (78,4) rectangle (118,37);
  \node[hd, anchor=north west] at (78.5,36.5) {\iflangja{ノード 9（管理ノード）}{node 9 (manager)}};
  \node[lab, anchor=north west, text=ln-muted] at (78.5,32) {\iflangja{管理するページのメタデータ}{metadata of its pages}};
  \foreach \cx/\cw/\ct in {86.5/13/{\iflangja{ページ}{page}}, 99/12/{\iflangja{所有者}{owner}}, 110.5/11/{\iflangja{複製}{replicas}}} {
    \node[cellc, minimum width=\cw mm, fill=ln-box] at (\cx,24) {\ct};
  }
  \foreach \ry/\pa/\pb/\pc in {19/\ttfamily 050012A3/2/{4, 7}, 14/\ttfamily 050012A4/9/---} {
    \node[cellc, minimum width=13mm, fill=white] at (86.5,\ry) {\pa};
    \node[cellc, minimum width=12mm, fill=white] at (99,\ry) {\pb};
    \node[cellc, minimum width=11mm, fill=white] at (110.5,\ry) {\pc};
  }
  \node[lab, text=ln-muted] at (98,8.5) {$\vdots$};
  \draw[->] (70,21) -- (78,21);
\end{tikzpicture}
   (width ~118 mm)
\begin{tikzpicture}[x=1mm, y=1mm, font=\scriptsize\sffamily,
  >={Stealth[length=1.4mm]},
  fld/.style={draw, minimum height=6mm, inner sep=0pt, font=\scriptsize\ttfamily},
  cellc/.style={draw, minimum height=5mm, inner sep=0pt},
  mach/.style={draw, rounded corners=0.8mm, fill=white},
  lab/.style={font=\scriptsize\sffamily, inner sep=0.8pt},
  hd/.style={font=\scriptsize\sffamily\bfseries, inner sep=0.8pt}]
  % page ID
  \node[hd, anchor=south west] at (0,37) {\iflangja{ページ ID}{page ID}};
  \node[fld, fill=ln-accent!20, minimum width=12mm] (nid) at (6,33) {05};
  \node[fld, fill=white, minimum width=22mm] at (23,33) {0012A3};
  \node[lab, anchor=north, text=ln-muted] at (6,29.6) {\iflangja{ノード ID}{node ID}};
  \node[lab, anchor=north, text=ln-muted] at (23,29.6) {PFN};
  % global mapping table
  \node[hd, anchor=south] at (58,37) {\iflangja{グローバルマップ}{global mapping table}};
  \node[lab, anchor=south, text=ln-muted] at (58,34) {\iflangja{（全ノードに複製）}{(replicated on every node)}};
  \node[lab, text=ln-muted] at (51,30.5) {\iflangja{索引}{index}};
  \node[lab, text=ln-muted] at (63,30.5) {\iflangja{管理ノード}{manager}};
  \foreach \ry/\ri/\rm in {26/04/4, 21/05/9, 16/06/6} {
    \node[cellc, minimum width=10mm, fill=white] at (51,\ry) {\ttfamily\ri};
    \node[cellc, minimum width=14mm, fill=white] at (63,\ry) {\rm};
  }
  \node[cellc, minimum width=14mm, fill=ln-accent!20] at (63,21) {9};
  \node[lab, text=ln-muted] at (51,11) {$\vdots$};
  \node[lab, text=ln-muted] at (63,11) {$\vdots$};
  \draw[->] (nid.south) ++(0,-3.6) -- (6,21) -- (46,21);
  % manager node
  \draw[mach] (78,4) rectangle (118,37);
  \node[hd, anchor=north west] at (78.5,36.5) {\iflangja{ノード 9（管理ノード）}{node 9 (manager)}};
  \node[lab, anchor=north west, text=ln-muted] at (78.5,32) {\iflangja{管理するページのメタデータ}{metadata of its pages}};
  \foreach \cx/\cw/\ct in {86.5/13/{\iflangja{ページ}{page}}, 99/12/{\iflangja{所有者}{owner}}, 110.5/11/{\iflangja{複製}{replicas}}} {
    \node[cellc, minimum width=\cw mm, fill=ln-box] at (\cx,24) {\ct};
  }
  \foreach \ry/\pa/\pb/\pc in {19/\ttfamily 050012A3/2/{4, 7}, 14/\ttfamily 050012A4/9/---} {
    \node[cellc, minimum width=13mm, fill=white] at (86.5,\ry) {\pa};
    \node[cellc, minimum width=12mm, fill=white] at (99,\ry) {\pb};
    \node[cellc, minimum width=11mm, fill=white] at (110.5,\ry) {\pc};
  }
  \node[lab, text=ln-muted] at (98,8.5) {$\vdots$};
  \draw[->] (70,21) -- (78,21);
\end{tikzpicture}

  \caption{Indexing distributed state.  The node-ID field of the page ID
    indexes the global mapping table, which names the manager of the page;
    the manager keeps the metadata of its pages.  Entry \texttt{05} has
    been changed to node~9, which therefore manages the pages whose node-ID
    field is \texttt{05}, although their frames reside on node~5.}
  \label{fig:l15-indexing}
\end{figure}

\begin{example}
  A DSM identifies each page by a 32-bit address: an 8-bit node ID followed
  by a 24-bit PFN, so that each of up to 256 nodes can contribute up to
  $2^{24}$ frames of \SI{4}{\kibi\byte}, \SI{64}{\gibi\byte} of memory.
  Page \texttt{0x050012A3} is frame \texttt{0x0012A3} of node~5, which
  manages it initially.  The mapping table is indexed by the node-ID field
  and maps every index $i$ to node~$i$ to begin with.  When node~5 becomes
  overloaded, the management of its pages is handed to node~9 by setting
  entry~5 to 9 on all nodes and moving the page metadata from node~5 to
  node~9 (\cref{fig:l15-indexing}); the frames themselves stay on node~5,
  and no page ID changes.  Had the manager been taken directly from the
  node ID field, every one of these pages would have needed a new ID, and
  every reference to it would have had to be updated.
\end{example}

\section{Implementing DSM}
\label{sec:l15-implementing}

A DSM must intercept accesses to the shared memory for two
reasons.\source{P4L3, implementing DSM; Nitzberg and Lo; Tanenbaum and Bos
ch.~8.}  When a process accesses memory that is not present on its node, the DSM
must send messages to the node that holds it, requesting access.  When a
process updates memory of which other nodes hold replicas, the DSM must
detect the update and perform the consistency operations it requires.  At
the same time, accesses to memory that is present locally and needs no
coordination must involve neither the DSM nor the operating system: an
intervention by software in every memory access would forfeit the
performance that locality is meant to provide.

A page-based DSM resolves this conflict with the hardware support of the
MMU (\cref{fig:l15-mmu}).  The MMU traps into the operating system with a
page fault\index{page fault} (Lesson~8) whenever a virtual address has no
valid mapping, or the mapping does not permit the kind of access attempted;
every other access completes in hardware.\caution{Such a page fault is not
served from swap space: the DSM fetches the page from the memory of another
node.}  The DSM sets up the mappings so that exactly the accesses it must
intercept fault.
\begin{itemize}
  \item A page that is not present on the node has no valid mapping.  An
    access to it traps, and the operating system passes the fault to the
    DSM, which requests the page from its owner and installs a mapping for
    the received copy before it resumes the process: read-only for a read;
    for a write, the DSM performs the consistency operations below as well
    and maps the page read-write, so that a write to an absent page costs
    one fault and not two.
  \item A page of which the node holds a replica is mapped read-only
    through its protection bits\index{protection bits} (Lesson~8).  Reads
    proceed without traps.  A write traps, and the DSM performs the
    consistency operations---it invalidates the other replicas and takes
    over ownership---before it enables writing and resumes the
    process.\sidenote{IVY \cite{li1989}, the first page-based DSM, works in
    this way and provides sequential consistency under MRSW access;
    TreadMarks \cite{keleher1994} later implemented a DSM entirely in user
    space, on unmodified UNIX systems, with a weaker model.}
\end{itemize}
Other information recorded by the MMU is useful as well.  The dirty
bit\index{dirty bit} of a page shows whether the page has been modified, and
hence whether modifications must be propagated; the access
bit\index{access bit} shows whether a local replica has been used recently,
which helps to decide which replicas to evict when memory runs short, as
in page replacement (Lesson~8).  An object-based DSM intercepts accesses in
the language runtime instead, which mediates the program's operations on
shared objects.\sidenote{These MMU operations are not free.  On the systems
of 1991 Appel and Li \cite{appel1991} measured some
\SI{1200}{\micro\second} of fault and re-protection overhead on an average
one, plus a TLB shootdown on the other CPUs that hold the mapping.}

\begin{figure}[htbp]
  \centering
  % Implementing a page-based DSM with the MMU: accesses permitted by the
% current mapping complete in hardware; the others fault, and the DSM fetches
% the page, performs consistency operations, or both, before the access is
% repeated.
% Usage: % Implementing a page-based DSM with the MMU: accesses permitted by the
% current mapping complete in hardware; the others fault, and the DSM fetches
% the page, performs consistency operations, or both, before the access is
% repeated.
% Usage: % Implementing a page-based DSM with the MMU: accesses permitted by the
% current mapping complete in hardware; the others fault, and the DSM fetches
% the page, performs consistency operations, or both, before the access is
% repeated.
% Usage: \input{figures/l15-mmu}   (width ~112 mm)
\begin{tikzpicture}[x=1mm, y=1mm, font=\scriptsize\sffamily,
  >={Stealth[length=1.5mm]},
  box/.style={draw, fill=white, align=center, inner sep=1.2pt, minimum height=8mm, text width=30mm},
  dec/.style={draw, diamond, aspect=2.2, fill=ln-box, align=center, inner sep=0pt},
  lab/.style={font=\scriptsize\sffamily, inner sep=0.8pt}]
  \node[box] (acc) at (17,50) {\iflangja{プロセスが共有アドレス\\にアクセス}{process accesses a\\shared address}};
  \node[dec] (mmu) at (17,33) {\iflangja{MMU：アクセス\\は許可？}{MMU: access\\permitted?}};
  \node[box, fill=ln-tip!12] (ok) at (17,14) {\iflangja{ハードウェアで完了\\（トラップなし）}{access completes in\\hardware (no trap)}};
  \node[box, fill=ln-source!12, text width=34mm] (trap) at (66,33) {\iflangja{ページフォールト：\\OS へトラップ，\\OS は DSM に渡す}{page fault: trap to the OS,\\which hands it to the DSM}};
  \node[box] (fetch) at (50,10) {\iflangja{所有者にページを要求：\\読みなら読み出し専用\\でマップ}{request the page from its\\owner; map it read-only\\for a read}};
  \node[box, text width=32mm] (coord) at (93,10) {\iflangja{一貫性操作：他の複製を\\無効化，読み書き可でマップ}{consistency operations:\\invalidate other replicas;\\map it read-write}};
  \draw[->] (acc) -- (mmu);
  \draw[->] (mmu) -- node[lab, right] {\iflangja{はい}{yes}} (ok);
  \draw[->] (mmu) -- node[lab, above] {\iflangja{いいえ}{no}} (trap);
  \draw[->] (trap.south) -- ++(0,-5) -| (fetch.north);
  \draw[->] (trap.south) -- ++(0,-5) -| (coord.north);
  \node[lab, anchor=east, text=ln-muted, align=right] at (49.2,20) {\iflangja{ローカルに\\ない}{no local\\copy}};
  \draw[->] (fetch.east) -- node[lab, above, inner sep=0.6pt] {\iflangja{書き}{write}} (coord.west);
  \node[lab, anchor=west, text=ln-muted] at (50.6,4.2) {\iflangja{読み}{read}};
  \node[lab, anchor=west, text=ln-muted, align=left] at (93.8,20) {\iflangja{読み出し専用\\への書き込み}{write to a\\read-only copy}};
  \draw[->] (fetch.south) -- ++(0,-5) -| (ok.south);
  \draw (coord.south) |- ++(-40,-5);
  \node[lab, anchor=north, text=ln-muted] at (73,-0.4) {\iflangja{プロセスを再開：アクセスを再実行}{resume the process: the access is repeated}};
\end{tikzpicture}
   (width ~112 mm)
\begin{tikzpicture}[x=1mm, y=1mm, font=\scriptsize\sffamily,
  >={Stealth[length=1.5mm]},
  box/.style={draw, fill=white, align=center, inner sep=1.2pt, minimum height=8mm, text width=30mm},
  dec/.style={draw, diamond, aspect=2.2, fill=ln-box, align=center, inner sep=0pt},
  lab/.style={font=\scriptsize\sffamily, inner sep=0.8pt}]
  \node[box] (acc) at (17,50) {\iflangja{プロセスが共有アドレス\\にアクセス}{process accesses a\\shared address}};
  \node[dec] (mmu) at (17,33) {\iflangja{MMU：アクセス\\は許可？}{MMU: access\\permitted?}};
  \node[box, fill=ln-tip!12] (ok) at (17,14) {\iflangja{ハードウェアで完了\\（トラップなし）}{access completes in\\hardware (no trap)}};
  \node[box, fill=ln-source!12, text width=34mm] (trap) at (66,33) {\iflangja{ページフォールト：\\OS へトラップ，\\OS は DSM に渡す}{page fault: trap to the OS,\\which hands it to the DSM}};
  \node[box] (fetch) at (50,10) {\iflangja{所有者にページを要求：\\読みなら読み出し専用\\でマップ}{request the page from its\\owner; map it read-only\\for a read}};
  \node[box, text width=32mm] (coord) at (93,10) {\iflangja{一貫性操作：他の複製を\\無効化，読み書き可でマップ}{consistency operations:\\invalidate other replicas;\\map it read-write}};
  \draw[->] (acc) -- (mmu);
  \draw[->] (mmu) -- node[lab, right] {\iflangja{はい}{yes}} (ok);
  \draw[->] (mmu) -- node[lab, above] {\iflangja{いいえ}{no}} (trap);
  \draw[->] (trap.south) -- ++(0,-5) -| (fetch.north);
  \draw[->] (trap.south) -- ++(0,-5) -| (coord.north);
  \node[lab, anchor=east, text=ln-muted, align=right] at (49.2,20) {\iflangja{ローカルに\\ない}{no local\\copy}};
  \draw[->] (fetch.east) -- node[lab, above, inner sep=0.6pt] {\iflangja{書き}{write}} (coord.west);
  \node[lab, anchor=west, text=ln-muted] at (50.6,4.2) {\iflangja{読み}{read}};
  \node[lab, anchor=west, text=ln-muted, align=left] at (93.8,20) {\iflangja{読み出し専用\\への書き込み}{write to a\\read-only copy}};
  \draw[->] (fetch.south) -- ++(0,-5) -| (ok.south);
  \draw (coord.south) |- ++(-40,-5);
  \node[lab, anchor=north, text=ln-muted] at (73,-0.4) {\iflangja{プロセスを再開：アクセスを再実行}{resume the process: the access is repeated}};
\end{tikzpicture}
   (width ~112 mm)
\begin{tikzpicture}[x=1mm, y=1mm, font=\scriptsize\sffamily,
  >={Stealth[length=1.5mm]},
  box/.style={draw, fill=white, align=center, inner sep=1.2pt, minimum height=8mm, text width=30mm},
  dec/.style={draw, diamond, aspect=2.2, fill=ln-box, align=center, inner sep=0pt},
  lab/.style={font=\scriptsize\sffamily, inner sep=0.8pt}]
  \node[box] (acc) at (17,50) {\iflangja{プロセスが共有アドレス\\にアクセス}{process accesses a\\shared address}};
  \node[dec] (mmu) at (17,33) {\iflangja{MMU：アクセス\\は許可？}{MMU: access\\permitted?}};
  \node[box, fill=ln-tip!12] (ok) at (17,14) {\iflangja{ハードウェアで完了\\（トラップなし）}{access completes in\\hardware (no trap)}};
  \node[box, fill=ln-source!12, text width=34mm] (trap) at (66,33) {\iflangja{ページフォールト：\\OS へトラップ，\\OS は DSM に渡す}{page fault: trap to the OS,\\which hands it to the DSM}};
  \node[box] (fetch) at (50,10) {\iflangja{所有者にページを要求：\\読みなら読み出し専用\\でマップ}{request the page from its\\owner; map it read-only\\for a read}};
  \node[box, text width=32mm] (coord) at (93,10) {\iflangja{一貫性操作：他の複製を\\無効化，読み書き可でマップ}{consistency operations:\\invalidate other replicas;\\map it read-write}};
  \draw[->] (acc) -- (mmu);
  \draw[->] (mmu) -- node[lab, right] {\iflangja{はい}{yes}} (ok);
  \draw[->] (mmu) -- node[lab, above] {\iflangja{いいえ}{no}} (trap);
  \draw[->] (trap.south) -- ++(0,-5) -| (fetch.north);
  \draw[->] (trap.south) -- ++(0,-5) -| (coord.north);
  \node[lab, anchor=east, text=ln-muted, align=right] at (49.2,20) {\iflangja{ローカルに\\ない}{no local\\copy}};
  \draw[->] (fetch.east) -- node[lab, above, inner sep=0.6pt] {\iflangja{書き}{write}} (coord.west);
  \node[lab, anchor=west, text=ln-muted] at (50.6,4.2) {\iflangja{読み}{read}};
  \node[lab, anchor=west, text=ln-muted, align=left] at (93.8,20) {\iflangja{読み出し専用\\への書き込み}{write to a\\read-only copy}};
  \draw[->] (fetch.south) -- ++(0,-5) -| (ok.south);
  \draw (coord.south) |- ++(-40,-5);
  \node[lab, anchor=north, text=ln-muted] at (73,-0.4) {\iflangja{プロセスを再開：アクセスを再実行}{resume the process: the access is repeated}};
\end{tikzpicture}

  \caption{Intercepting accesses in a page-based DSM.  Accesses permitted by
    the current mapping never leave the hardware; the others fault, and the
    DSM fetches the page or performs consistency operations before the
    access is repeated.}
  \label{fig:l15-mmu}
\end{figure}

\begin{example}
  Page $p$ initially resides only on node~1, its owner, which maps it
  read-write.  \Cref{tab:l15-mmu} traces accesses by three nodes under MRSW
  access with invalidation on write.\tip{Per node the mapping of a page is
  a three-state machine---none, RO, RW---and only the transitions fault: a
  read takes none to RO, a write RO to RW; another node's access
  downgrades it to RO or to none.}  When node~2 obtains a replica in
  step~1, node~1's mapping is downgraded to read-only, so that a later write
  by node~1 would fault and invalidate the replica.  Of the seven accesses,
  four fault; the reads in steps~2 and~3 and the write in step~5 complete
  in hardware, because the mappings already permit them.
\end{example}

\begin{table}[htbp]
  \centering
  \caption{Accesses to page $p$ and the faults they cause.  The last column
    shows the mappings of $p$ on nodes 1/2/3 after each step (RW
    read-write, RO read-only, --- none).}
  \label{tab:l15-mmu}
  \small
  \begin{tabular}{@{}rll>{\raggedright\arraybackslash}p{0.34\linewidth}l@{}}
    \toprule
    Step & Access & Fault & DSM action & Maps \\
    \midrule
    0 & (initial) & --- & --- & RW/---/--- \\
    1 & node 2 reads & page fault & fetch $p$ from owner node~1; downgrade
      node~1 to RO & RO/RO/--- \\
    2 & node 2 reads & none & --- & RO/RO/--- \\
    3 & node 1 reads & none & --- & RO/RO/--- \\
    4 & node 2 writes & protection & invalidate node~1's copy; node~2
      becomes owner & ---/RW/--- \\
    5 & node 2 writes & none & --- & ---/RW/--- \\
    6 & node 3 reads & page fault & fetch $p$ from owner node~2; downgrade
      node~2 to RO & ---/RO/RO \\
    7 & node 2 writes & protection & invalidate node~3's copy & ---/RW/--- \\
    \bottomrule
  \end{tabular}
\end{table}

\begin{keypoint}
  A page-based DSM places the shared pages in memory contributed by all
  nodes, finds each page's manager through a replicated global map, and
  uses the MMU to intercept exactly the accesses that need it: accesses to
  pages that are not present, and writes to replicated pages.  All other
  accesses proceed at hardware speed.
\end{keypoint}

\Needspace{6\baselineskip}
\section{Consistency models}
\label{sec:l15-models}

The mechanisms of \cref{sec:l15-management} can be triggered at different
times, and the resulting guarantees differ.\source{P4L3, consistency
model; strict, sequential, causal and weak consistency; Tanenbaum and van
Steen \cite{tanenbaum2007} ch.~7.}  For a
distributed file system these guarantees are the file sharing semantics of
Lesson~14; for memory, and for shared state in general, they are stated by
a consistency model.

\begin{definition}[Consistency model]
  A \term{consistency model}[一貫性モデル] is an agreement between memory,
  or any shared state, and the software layers above it: the memory
  guarantees to behave correctly---with respect to the ordering of
  accesses and to the propagation and visibility of updates---provided
  that the software follows specific rules.
\end{definition}

File sharing semantics are thus consistency models for files.  To describe
consistency models for memory we write $R_{m}(v)$ for a read
that returns the value $v$ from memory location $m$, and $W_{m}(v)$ for a
write of $v$ to $m$.  Operations are drawn on one timeline per process,
with time increasing to the right; all locations initially hold
0.\margin{In the timelines $m_a$ and $m_b$ are two locations of the shared
address space, and $t_1 < t_2 < \dots$ are instants of real time as an
outside observer sees them.}

\subsection{Strict consistency}

\begin{definition}[Strict consistency]
  Under \term{strict consistency}[厳密一貫性] every update is visible
  everywhere immediately: every read, on any node, returns the value of the
  write to that location that is the most recent in real time, and all
  nodes observe the updates in the order in which they were made.
\end{definition}

Strict consistency is what programmers intuitively expect of memory, but it
is a theoretical model.  Even on a single shared memory multiprocessor,
concurrent updates on different CPUs take effect in no guaranteed order
unless the software adds locking and synchronization.  In a distributed
system, updates travel over a network with non-zero latency, messages may
be reordered or lost, and no global clock defines which of two writes on
different nodes is the more recent; strict consistency cannot be
guaranteed.

\subsection{Sequential consistency}

\begin{definition}[Sequential consistency]
  Under \term{sequential consistency}[逐次一貫性] the result of every
  execution is the same as if the operations of all processes were executed
  in some sequential order, in which the operations of each process appear
  in the order in which the process issued them, every read returns the
  most recent write to its location, and all processes observe this same
  order \cite{lamport1979}.
\end{definition}

Updates from different processes may thus be interleaved arbitrarily, and
the interleaving need not match the real-time order, but all processes
must see the same interleaving.\caution{A sequentially consistent execution
need not agree with real time: a read may return a value older than the
write that most recently preceded it in real time, provided that every
process observes the same order.}

\subsection{Causal consistency}

A write \emph{causally depends} on an earlier write if the process that
issued it had observed the earlier write before---for instance by reading
the value it wrote---directly or through a chain of such dependencies.
Writes of which neither depends on the other are \emph{concurrent}.

\begin{definition}[Causal consistency]
  Under \term{causal consistency}[因果一貫性] all processes observe causally
  related writes in their causal order, and the operations of each process
  in the order in which the process issued them; concurrent writes may be
  observed in different orders by different processes.
\end{definition}

To provide causal consistency, the DSM must detect possible causal
relationships, that is, record which updates a process had observed before
each of its writes.\margin{Vector timestamps are a common way to record
these dependencies (Tanenbaum and van Steen ch.~6).}  In exchange, it needs
to coordinate only the order of causally related updates.

\subsection{Weak consistency}

Weak consistency introduces \emph{synchronization points}: operations, such
as \texttt{sync}, that the DSM makes available to the software above it.

\begin{definition}[Weak consistency]
  Under \term{weak consistency}[弱一貫性] every update made before a
  synchronization point is visible to all nodes after it; between
  synchronization points no guarantee is given about which updates a node
  observes.
\end{definition}

A DSM may offer a single synchronization operation for the entire shared
state, or separate operations for subsets of the state, such as individual
pages or objects, so that an application synchronizes only the state it
needs.  It may also split the operation in two: an \emph{acquire},
performed before a process uses shared state, which makes the updates of
the other processes visible to it, and a \emph{release}, performed after it
has finished updating, which propagates its own updates.\margin{Session
semantics (Lesson~14) is the file counterpart of weak consistency:
\texttt{open} and \texttt{close} are the synchronization points.}  Weak
consistency limits data movement and coordination to the synchronization
points, at which the updates made since the previous point can be
propagated together.  In exchange, the DSM must maintain extra state
to record and aggregate these updates, and between synchronization points
processes may observe stale data.

\subsection{Comparing the models}

Strict, sequential and causal consistency form a hierarchy: every strictly
consistent execution is sequentially consistent, and every sequentially
consistent execution is causally consistent.  A weaker model admits more
executions and therefore requires less coordination, but it gives the
applications fewer guarantees.\caution{Real hardware is not sequentially
consistent either: x86 lets a load pass an earlier store to another location
(total store order), ARM and POWER reorder more, and fences take the part of
synchronization points (Intel SDM vol.~3A, memory ordering).}

\begin{figure}[htbp]
  \centering
  % Operations for the consistency-model example: P1 and P2 write m_a and m_b,
% P3 and P4 read both locations in opposite orders, a synchronization point
% follows at t5, and P4 reads m_a again at t6.  P2's read of m_a (grey) is
% present only in the variant with a causal dependency.
% Usage: % Operations for the consistency-model example: P1 and P2 write m_a and m_b,
% P3 and P4 read both locations in opposite orders, a synchronization point
% follows at t5, and P4 reads m_a again at t6.  P2's read of m_a (grey) is
% present only in the variant with a causal dependency.
% Usage: % Operations for the consistency-model example: P1 and P2 write m_a and m_b,
% P3 and P4 read both locations in opposite orders, a synchronization point
% follows at t5, and P4 reads m_a again at t6.  P2's read of m_a (grey) is
% present only in the variant with a causal dependency.
% Usage: \input{figures/l15-models}   (width ~112 mm)
\begin{tikzpicture}[x=1mm, y=1mm, font=\scriptsize\sffamily,
  row/.style={font=\scriptsize\sffamily, anchor=east},
  lab/.style={font=\scriptsize, inner sep=0.8pt},
  ev/.style={circle, fill=black, inner sep=0pt, minimum size=1.3mm}]
  % t_i -> x = 8 + 16 i (mm); rows: P1 27, P2 18, P3 9, P4 0
  \foreach \pn/\py in {1/27,2/18,3/9,4/0} {
    \node[row] at (14,\py) {P\pn};
    \draw[ln-rule] (16,\py) -- (112,\py);
  }
  % synchronization point
  \draw[densely dashed, ln-accent, thick] (88,-3) -- (88,31);
  \node[lab, anchor=south, text=ln-accent, font=\scriptsize\sffamily] at (88,31) {sync};
  % P1, P2 writes
  \node[ev] at (24,27) {};  \node[lab, anchor=south] at (24,28) {$W_{m_a}(7)$};
  \node[ev] at (40,18) {};  \node[lab, anchor=south] at (40,19) {$W_{m_b}(9)$};
  \node[ev, fill=ln-muted] at (31,18) {};
  \node[lab, anchor=north, text=ln-muted] at (29,17) {($R_{m_a}(7)$)};
  % P3 reads
  \node[ev] at (56,9) {};  \node[lab, anchor=south] at (56,10) {$R_{m_a}(?)$};
  \node[ev] at (72,9) {};  \node[lab, anchor=south] at (72,10) {$R_{m_b}(?)$};
  % P4 reads
  \node[ev] at (56,0) {};  \node[lab, anchor=south] at (56,1) {$R_{m_b}(?)$};
  \node[ev] at (72,0) {};  \node[lab, anchor=south] at (72,1) {$R_{m_a}(?)$};
  \node[ev] at (104,0) {}; \node[lab, anchor=south] at (104,1) {$R_{m_a}(?)$};
  % time axis
  \draw[->] (16,-6) -- (113,-6);
  \foreach \ti in {1,...,6} {
    \draw ({8+16*\ti},-5.2) -- ({8+16*\ti},-6.8);
    \node[lab, anchor=north] at ({8+16*\ti},-7) {$t_{\ti}$};
  }
\end{tikzpicture}
   (width ~112 mm)
\begin{tikzpicture}[x=1mm, y=1mm, font=\scriptsize\sffamily,
  row/.style={font=\scriptsize\sffamily, anchor=east},
  lab/.style={font=\scriptsize, inner sep=0.8pt},
  ev/.style={circle, fill=black, inner sep=0pt, minimum size=1.3mm}]
  % t_i -> x = 8 + 16 i (mm); rows: P1 27, P2 18, P3 9, P4 0
  \foreach \pn/\py in {1/27,2/18,3/9,4/0} {
    \node[row] at (14,\py) {P\pn};
    \draw[ln-rule] (16,\py) -- (112,\py);
  }
  % synchronization point
  \draw[densely dashed, ln-accent, thick] (88,-3) -- (88,31);
  \node[lab, anchor=south, text=ln-accent, font=\scriptsize\sffamily] at (88,31) {sync};
  % P1, P2 writes
  \node[ev] at (24,27) {};  \node[lab, anchor=south] at (24,28) {$W_{m_a}(7)$};
  \node[ev] at (40,18) {};  \node[lab, anchor=south] at (40,19) {$W_{m_b}(9)$};
  \node[ev, fill=ln-muted] at (31,18) {};
  \node[lab, anchor=north, text=ln-muted] at (29,17) {($R_{m_a}(7)$)};
  % P3 reads
  \node[ev] at (56,9) {};  \node[lab, anchor=south] at (56,10) {$R_{m_a}(?)$};
  \node[ev] at (72,9) {};  \node[lab, anchor=south] at (72,10) {$R_{m_b}(?)$};
  % P4 reads
  \node[ev] at (56,0) {};  \node[lab, anchor=south] at (56,1) {$R_{m_b}(?)$};
  \node[ev] at (72,0) {};  \node[lab, anchor=south] at (72,1) {$R_{m_a}(?)$};
  \node[ev] at (104,0) {}; \node[lab, anchor=south] at (104,1) {$R_{m_a}(?)$};
  % time axis
  \draw[->] (16,-6) -- (113,-6);
  \foreach \ti in {1,...,6} {
    \draw ({8+16*\ti},-5.2) -- ({8+16*\ti},-6.8);
    \node[lab, anchor=north] at ({8+16*\ti},-7) {$t_{\ti}$};
  }
\end{tikzpicture}
   (width ~112 mm)
\begin{tikzpicture}[x=1mm, y=1mm, font=\scriptsize\sffamily,
  row/.style={font=\scriptsize\sffamily, anchor=east},
  lab/.style={font=\scriptsize, inner sep=0.8pt},
  ev/.style={circle, fill=black, inner sep=0pt, minimum size=1.3mm}]
  % t_i -> x = 8 + 16 i (mm); rows: P1 27, P2 18, P3 9, P4 0
  \foreach \pn/\py in {1/27,2/18,3/9,4/0} {
    \node[row] at (14,\py) {P\pn};
    \draw[ln-rule] (16,\py) -- (112,\py);
  }
  % synchronization point
  \draw[densely dashed, ln-accent, thick] (88,-3) -- (88,31);
  \node[lab, anchor=south, text=ln-accent, font=\scriptsize\sffamily] at (88,31) {sync};
  % P1, P2 writes
  \node[ev] at (24,27) {};  \node[lab, anchor=south] at (24,28) {$W_{m_a}(7)$};
  \node[ev] at (40,18) {};  \node[lab, anchor=south] at (40,19) {$W_{m_b}(9)$};
  \node[ev, fill=ln-muted] at (31,18) {};
  \node[lab, anchor=north, text=ln-muted] at (29,17) {($R_{m_a}(7)$)};
  % P3 reads
  \node[ev] at (56,9) {};  \node[lab, anchor=south] at (56,10) {$R_{m_a}(?)$};
  \node[ev] at (72,9) {};  \node[lab, anchor=south] at (72,10) {$R_{m_b}(?)$};
  % P4 reads
  \node[ev] at (56,0) {};  \node[lab, anchor=south] at (56,1) {$R_{m_b}(?)$};
  \node[ev] at (72,0) {};  \node[lab, anchor=south] at (72,1) {$R_{m_a}(?)$};
  \node[ev] at (104,0) {}; \node[lab, anchor=south] at (104,1) {$R_{m_a}(?)$};
  % time axis
  \draw[->] (16,-6) -- (113,-6);
  \foreach \ti in {1,...,6} {
    \draw ({8+16*\ti},-5.2) -- ({8+16*\ti},-6.8);
    \node[lab, anchor=north] at ({8+16*\ti},-7) {$t_{\ti}$};
  }
\end{tikzpicture}

  \caption{Operations of \cref{ex:l15-models}.  P2's read of $m_a$ (grey)
    occurs only in the variant in which P2's write depends on P1's.  The
    dashed line is a synchronization point.}
  \label{fig:l15-models}
\end{figure}

\begin{example}
  \label{ex:l15-models}
  In \cref{fig:l15-models}, P1 writes 7 to $m_a$ at $t_1$ and P2 writes 9 to
  $m_b$ at $t_2$.  P3 and P4 read both locations at $t_3$ and $t_4$, in
  opposite orders.  \Cref{tab:l15-models} lists three possible outcomes of
  these reads and the models they satisfy.
  \begin{itemize}
    \item Outcome (a) is strictly consistent: every read returns the most
      recent write in real time.
    \item Outcome (b) is not: P4 reads 0 from $m_b$ at $t_3$, after 9 was
      written at $t_2$.  It is sequentially consistent, because the order
      $R_{m_a}(0)$ by P3, $R_{m_b}(0)$ by P4, $W_{m_a}(7)$, $W_{m_b}(9)$,
      $R_{m_b}(9)$ by P3, $R_{m_a}(7)$ by P4 explains every read and
      preserves the order of each process's operations.\tip{To test
      sequential consistency, look for one merged list of the operations
      that keeps each process's own order and explains every read.}
    \item Outcome (c) is not sequentially consistent: P3 observes
      $W_{m_a}(7)$ before $W_{m_b}(9)$, P4 observes them the other way
      round, and no single order explains both.  If the two writes are
      concurrent, (c) is causally consistent.  If P2 read 7 from $m_a$
      before writing $m_b$, however, $W_{m_b}(9)$ causally depends on
      $W_{m_a}(7)$, and P4, which reads 9 from $m_b$ and afterwards 0 from
      $m_a$, observes the effect before its cause.
  \end{itemize}
  With a synchronization point at $t_5$, all three outcomes are weakly
  consistent, because no guarantee applies before $t_5$; the read of $m_a$
  by P4 at $t_6$, however, must return 7.\caution{A model constrains the
  outcomes; it does not determine them.  A DSM that provides sequential
  consistency may produce (a) on one run and (b) on the next, and a weakly
  consistent one may produce any of the three.}
\end{example}

\begin{table}[htbp]
  \centering
  \caption{Outcomes of the reads at $t_3$ and $t_4$ in \cref{ex:l15-models}.
    The causal column gives the result for concurrent writes / for P2
    reading $m_a$ before writing $m_b$.}
  \label{tab:l15-models}
  \small
  \begin{tabular}{@{}lllcccc@{}}
    \toprule
    & P3 & P4 & Strict & Sequential & Causal & Weak \\
    \midrule
    (a) & $R_{m_a}(7)$, $R_{m_b}(9)$ & $R_{m_b}(9)$, $R_{m_a}(7)$ & yes & yes
      & yes / yes & yes \\
    (b) & $R_{m_a}(0)$, $R_{m_b}(9)$ & $R_{m_b}(0)$, $R_{m_a}(7)$ & no & yes
      & yes / yes & yes \\
    (c) & $R_{m_a}(7)$, $R_{m_b}(0)$ & $R_{m_b}(9)$, $R_{m_a}(0)$ & no & no
      & yes / no & yes \\
    \bottomrule
  \end{tabular}
\end{table}

\begin{keypoint}[Lesson summary]
  Distributed shared memory is a peer distributed application in which every
  node contributes memory, serves accesses to it, and takes part in
  consistency protocols, so that processes on many machines share one
  address space.  Its design is determined by the granularity of sharing
  (pages or objects, with the risk of false sharing), the access algorithm
  (SRSW, MRSW, MRMW), the choice between migration and replication to
  achieve locality, and consistency management by pushing invalidations or
  pulling modifications.  A page-based DSM finds each page's manager through
  a replicated mapping table and relies on MMU faults to intercept only the
  accesses that need coordination.  Consistency models state what processes
  observe: strict consistency is unattainable in practice; sequential
  consistency requires one common order of all operations; causal
  consistency orders only causally related writes; and weak consistency
  guarantees visibility only at synchronization points.
\end{keypoint}

\end{document}
